% macro conflict, first definition kept: \newcommand{\Hom}{\operatorname{Hom}}
% macro conflict, first definition kept: \newcommand{\Ob}{\operatorname{Ob}}
% macro conflict, first definition kept: \newcommand{\Ind}{\mathrm{Ind}}
% macro conflict, first definition kept: \newcommand{\card}{\mathrm{card}}
% macro conflict, first definition kept: \newcommand{\Fl}{\mathrm{Fl}}
% macro conflict, first definition kept: \newcommand{\Ob}{\mathrm{Ob}}
% macro conflict, first definition kept: \newcommand{\Filt}{\mathrm{Filt}}
% macro conflict, first definition kept: \newcommand{\Ind}{\mathrm{Ind}}
% macro conflict, first definition kept: \newcommand{\Ob}{\operatorname{Ob}}
% macro conflict, first definition kept: \newcommand{\Ob}{\operatorname{Ob}}
% macro conflict, first definition kept: \newcommand{\Fl}{\mathrm{Fl}}
% macro conflict, first definition kept: \newcommand{\Ob}{\operatorname{Ob}}
% macro conflict, first definition kept: \newcommand{\Hom}{\operatorname{Hom}}
% macro conflict, first definition kept: \newcommand{\Ob}{\operatorname{Ob}}
% macro conflict, first definition kept: \newcommand{\Hom}{\operatorname{Hom}}
% macro conflict, first definition kept: \newcommand{\uHom}{\underline{\operatorname{Hom}}}
% macro conflict, first definition kept: \newcommand{\Ob}{\operatorname{Ob}}
\PassOptionsToPackage{hyperfootnotes=false}{hyperref}
\documentclass[11pt]{article}
\usepackage{fontspec}
\usepackage[a4paper,left=24mm,right=32mm,top=23mm,bottom=25mm,marginparwidth=22mm,marginparsep=4mm]{geometry}
\usepackage{microtype}
\usepackage{xurl}
\usepackage{marginfix}
\usepackage[english,french]{babel}
\usepackage{amsmath}
\usepackage{amssymb}
\usepackage{amsthm}
\usepackage{mathrsfs}
\usepackage{tikz-cd}
\usepackage{enumitem}
\setlength{\emergencystretch}{2em}
\setlength{\marginparpush}{5pt}
\AtBeginDocument{%
\begin{otherlanguage*}{english}
\thispagestyle{empty}
\vspace*{16mm}
\begin{center}
{\large LES D\'ERIVATEURS}\par
\vspace{9mm}
{\LARGE\bfseries Chapter XVIII, part a}\par
\vspace{5mm}
{\Large Pseudo-topoi, tensor products,\par
and cardinal filtrations}\par
\vspace{12mm}
{\large Text with commentary}\par
\vspace{9mm}
Manuscript pages 50--138\par
Edition printed pages 1--43\par
\end{center}
\vspace{14mm}
\noindent\textbf{Source.} Alexander Grothendieck, \emph{Les D\'erivateurs},
electronic edition deciphered and transcribed by M. K\"unzer and edited by
M. K\"unzer, J. Malgoire, and G. Maltsiniotis; version 17 December 2025.
\par\medskip
\noindent\textbf{Commentary and audit} AI-produced by
GPT-6 Astra (OpenAI, ChatGPT Pro mode), checked and integrated by Claude (Anthropic).
\par\medskip
\noindent The French source retains the edition's mathematical slips.
The English commentary is distinguished by a vertical bar and explains
the repairs. Results marked [ours] identify independently developed
formulations or proofs; the marker makes no claim of priority.
\par\medskip
\noindent Marginal labels preserve the edition's printed and manuscript
page systems. The register of claims and the table of corrections accompany
the text.
\vfill
\end{otherlanguage*}
\clearpage
}

\newcommand{\printedpage}[1]{\if\relax\detokenize{#1}\relax\else\marginpar{\footnotesize p.~#1}\fi}
\newcommand{\msp}[1]{\marginpar{\footnotesize [#1]}}
\newcommand{\uHom}{\underline{\mathrm{Hom}}}
\newcommand{\Ens}{\mathrm{Ens}}
\newcommand{\Fais}{\mathrm{Fais}}
\newcommand{\pstop}{\mathrm{pstop}}
\newcommand{\cM}{\mathcal{M}}
\newcommand{\cN}{\mathcal{N}}
\newcommand{\cX}{\mathcal{X}}
\newcommand{\cY}{\mathcal{Y}}
\newcommand{\cF}{\mathcal{F}}
\newcommand{\Hom}{\mathrm{Hom}}
\newcommand{\Top}{\mathrm{Top}}
\newcommand{\Ob}{\mathrm{Ob}\,}
\newcommand{\defeq}{\overset{\text{déf}}{=}}
\newcommand{\bt}{\boxtimes}
\newcommand{\tbt}{\mathbin{\widetilde{\boxtimes}}}
\newcommand{\linj}{\lhook\joinrel\longrightarrow}
\newcommand{\pr}{\operatorname{pr}}
\newcommand{\card}{\operatorname{card}}
\newcommand{\Fl}{\operatorname{Fl}}
\newcommand{\Filt}{\operatorname{Filt}}
\newcommand{\Ind}{\operatorname{Ind}}
\newcommand{\hookxrightarrow}[2][]{\lhook\joinrel\xrightarrow[#1]{#2}}
\newcommand{\PF}{\mathrm{PF}}
\newcommand{\Aut}{\operatorname{Aut}}
\newcommand{\uFilt}{\underline{\mathrm{Filt}}}
\newcommand{\circled}[1]{\tikz[baseline=(c.base)]\node[draw,circle,inner sep=1pt](c){#1};}
\begin{document}
% ---- from der-ch18_p001-005.tex
\printedpage{} % no page number printed
% head: (none)

\begin{center}
{\Large\textsf{LES DÉRIVATEURS}}

\bigskip

\textbf{ALEXANDRE GROTHENDIECK}

\bigskip\bigskip

\textbf{Chapitre XVIII}

\medskip

\textbf{Catégories et ensembles accessibles}
\end{center}

\newpage

\printedpage{} % no page number printed
% head: (none)

Ce texte a été déchiffré et transcrit en \LaTeXe{} par M. Künzer. Il a été édité par M. Künzer, J. Malgoire, et G. Maltsiniotis. La transcription est aussi fidèle que possible au manuscrit. Seules quelques corrections mineures évidentes ont été effectuées. Pour les rares ajouts ou commentaires des éditeurs, ainsi que pour la numérotation originale des pages du manuscrit, les caractères de machine à écrire \texttt{[typewriter]} entre crochets sont utilisés. Quelques notes de bas de page ajoutées par les éditeurs sont signalées par «~N.~Éd.~». La numérotation des pages de ce chapitre commence à la page 50, car à l'origine ce chapitre était la suite du chapitre XIX qui se termine au milieu de la page 50.

Cette édition est provisoire. Les remarques, commentaires et corrections sont bienvenus. Envoyer un message à :

\textbf{georges.maltsiniotis@imj-prg.fr}

\hfill G. Maltsiniotis\hspace*{2cm}

\newpage

\printedpage{1}
% head: A. Grothendieck : Dérivateurs, Ch. XVIII, Édition des Universités de Montpellier II et Paris VII
% foot: Version 17 décembre 2025

\msp{page 50}

\section*{4\quad Paratopos et l'opération $\mathcal{M}\boxtimes\mathcal{N}$}

\subsection*{4.1\quad Pseudo-topos. Problème du produit.}

J'appelle \textit{catégorie pseudo-$\mathfrak{U}$-topos}, ou simplement \textit{pseudo-topos}, une $\mathfrak{U}$-catégorie $\mathcal{M}$ stable par $\mathfrak{U}$-limites inductives, et qui admet une \textit{petite} sous-catégorie $\mathcal{M}_0$ \textit{strictement génératrice} (i.e. génératrice par épimorphismes stricts). Alors $\mathcal{M}$ admet aussi des $\mathfrak{U}$-limites projectives (SGA 4 I 8.12.7). (Mais l'exemple de (Gr), cf. loc. cit. 8.12.9, montre que $\mathcal{M}^\circ$ n'est par forcément, pour autant, un pseudo-$\mathfrak{U}$-topos, faute à $\mathcal{M}$ d'admettre une petite sous-catégorie strictement \textit{co}génératrice. Mais le théorème 2.6 \texttt{[du chapitre XIX]} donne des conditions sur $\mathcal{M}$ (surtout du type exactitude) % [sic] par forcément -> pas forcément, misprint; also the parenthesis opened at "(Mais" is never closed
\msp{page 51}
pour que $\mathcal{M}$ admette assez d'injectifs -- ce qui impliquera alors l'existence d'une petite sous-catégorie strictement cogénératrice, donc que $\mathcal{M}^\circ$ est également un pseudo-topos (\footnote{C'est le cas notamment si $\mathcal{M}$ est un topos.}).

Je visualise les catégories pseudo-topos comme étant la catégorie des “faisceaux” sur un objet $\mathcal{X}$, visualisé comme une sorte d'“espace” que je distingue de la catégorie $\mathcal{M}=\Fais(\mathcal{X})$ des faisceaux sur $\mathcal{X}$. C'est $\mathcal{X}$ qui sera pour ceci le “pseudo-topos”. Les pseudo-topos forment une 2-catégorie, la variance dans $\mathcal{X}$ étant opposée de celle dans $\Fais(\mathcal{X})=\mathcal{M}$. De façon précise, si $\mathcal{X}$ correspond à $\mathcal{M}$, $\mathcal{Y}$ à $\mathcal{N}$, un morphisme $f$ de $\mathcal{X}$ dans $\mathcal{Y}$ sera donné par un foncteur
\[
f^* \;:\; \mathcal{N} \longrightarrow \mathcal{M}
\tag*{(4.1.1)}
\]
en sens inverse, commutant aux $\varinjlim$ inductives quelconques (\footnote{Dire ici que les objets de $\mathcal{M}$ sont visualisés comme des “ouverts” généralisés de $\mathcal{X}$. On a $\mathcal{M}\xrightarrow{\sim}\uHom^!(\mathcal{M}^\circ,\Ens)$, i.e. $\mathcal{M}$ s'interprète bel et bien comme les “faisceaux” sur cette catégorie $\mathcal{M}$ d'ouverts \ldots}). (\textbf{N.B.} si $\mathcal{X}$, $\mathcal{Y}$ sont même des topos, i.e. $\mathcal{M}$, $\mathcal{N}$ des catégories-topos, et si on veut que $f$ soit un morphisme de topos, on exigera de plus que $f^*$ soit de plus \textit{exact à gauche}. On dira qu'un \textit{morphisme de pseudo-topos est exact} si de plus $f^*$ est exact à gauche.)

\msp{page 52}
En vertu du dual de loc. cit. 8.12.7, cela équivaut au fait que $f^*$ admette un adjoint à droite $f_*$,
\[
f_* \;:\; \mathcal{M} \longrightarrow \mathcal{N}\;,
\tag*{(4.1.2)}
\]
lequel commute donc aux petites $\varprojlim$\,. La donnée de $f_*$ satisfaisant à cette propriété détermine donc $f^*$ à isomorphisme unique près. Mais en général, il ne suffit pas de se donner un tel $f_*$, pour pouvoir affirmer qu'il corresponde à un $f^*$, i.e. qu'il admette un adjoint à droite\,; pour ceci, il faudrait qu'on sache que les foncteurs $\mathcal{M}\longrightarrow\Ens$ qui commutent aux petites $\varprojlim$ sont représentables, ce qui sera le cas si $\mathcal{M}$ admet une petite sous-catégorie \textit{co}génératrice, i.e. si $\mathcal{M}^\circ$ est également une catégorie pseudo-topos. % [sic] qu'il admette un adjoint à droite -> adjoint à gauche, f^* is the left adjoint of f_*

\newpage

\printedpage{2}
% head: A. Grothendieck : Dérivateurs, Ch. XVIII, Édition des Universités de Montpellier II et Paris VII
% foot: Version 17 décembre 2025

Il est sans doute judicieux de procéder comme dans le cas des topos, et d'appeler \textit{morphisme de pseudo-topos} $f:\mathcal{X}\longrightarrow\mathcal{Y}$ un quadruple
\[
(f^*,f_*,\alpha,\beta)
\tag*{(4.1.3)}
\]
\msp{page 53}
de deux foncteurs adjoints $(f^*,f_*)$, et des morphismes d'adjonction. Quand $f^*$ est une équivalence, ce qui équivaut à $f_*$ une équivalence, ou $\alpha$ et $\beta$ des isomorphismes, on dira qu'on a une \textit{équivalence de pseudo-topos}, et à toutes fins pratiques, on pourra alors regarder $\mathcal{X}$ et $\mathcal{Y}$ comme des objets essentiellement identiques.

Pour deux pseudo-topos $\mathcal{X}$, $\mathcal{Y}$, j'ai envie d'introduire (comme dans le cas des topos) un pseudo-topos produit
\[
\mathcal{X}\times\mathcal{Y}\;,
\tag*{(4.1.4)}
\]
qui ait les propriétés d'un 2-produit dans la 2-catégorie des pseudo-topos (\footnote{Il apparaît finalement que ce qu'on construit n'est \textit{pas} un 2-produit dans la 2-catégorie envisagée, donc il vaut mieux le noter par un signe tel que $\mathcal{X}\boxtimes\mathcal{Y}$. Mais si $\mathcal{X}$, $\mathcal{Y}$ sont deux topos, $\mathcal{X}\boxtimes\mathcal{Y}$ est un 2-produit dans la 2-catégorie des topos.}). Je devrais préciser que par définition
\[
\begin{aligned}
\uHom_{\pstop}(\mathcal{X},\mathcal{Y}) \;&\simeq\; \uHom_!(\mathcal{F}(\mathcal{Y}),\mathcal{F}(\mathcal{X}))^\circ\\
&\simeq\; \uHom^{!'}(\mathcal{F}(\mathcal{X}),\mathcal{F}(\mathcal{Y})),
\end{aligned}
\tag*{(4.1.5)}
\]
où l'indice ! (resp. l'exposant !) indique la sous-catégorie pleine dans la catégorie $\uHom$ envisagée formée des foncteurs qui commutent aux petites $\varinjlim$ (resp. aux petites
\msp{page 54}
$\varprojlim$\,), et l'accent dans le $\uHom^{!'}$ dans le dernier membre indique qu'on exige de plus que le foncteur en question ait un adjoint à gauche.

Soient donc
\[
\mathcal{M} \;=\; \mathcal{F}(\mathcal{X})\;, \qquad \mathcal{N} \;=\; \mathcal{F}(\mathcal{Y})\;,
\tag*{(4.1.6)}
\]
et soit $\mathcal{Z}$ un troisième pseudo-topos, décrit par
\[
\mathcal{P} \;=\; \mathcal{F}(\mathcal{Z})\;.
\tag*{(4.1.7)}
\]
On veut expliciter, à équivalence près, la catégorie produit $\uHom_{\pstop}(\mathcal{Z},\mathcal{X})\times\uHom_{\pstop}(\mathcal{Z},\mathcal{Y})$, i.e.
\[
\uHom_!(\mathcal{M},\mathcal{P})^\circ \;\times\; \uHom_!(\mathcal{N},\mathcal{Q})^\circ.
\tag*{(4.1.8)}
\] % [sic] \mathcal{Q} -> \mathcal{P}, \mathcal{Q} is not defined; \mathcal{Z} is described by \mathcal{P} (4.1.7)
Je vais introduire pour cela une catégorie
\printedpage{3}
% head: A. Grothendieck : Dérivateurs, Ch. XVIII, Édition des Universités de Montpellier II et Paris VII
% foot: Version 17 décembre 2025
\msp{page 55}
\[
\mathcal{M}\boxtimes\mathcal{N}
\tag*{(4.1.9)}
\]
par la formule
\[
\mathcal{M}\boxtimes\mathcal{N} \;\simeq\; \uHom^{!!}(\mathcal{M}^\circ\times\mathcal{N}^\circ,\Ens)\;,
\tag*{(4.1.10)}
\]
où le signe $\uHom^{!!}$ désigne la sous-catégorie pleine de la catégorie $\uHom$, formée des bi-contrafoncteurs % doubt resolved: retain the diplomatic bi-contrafoncteurs reading; see corrections.csv, PDF leaf 5
\[
F \;:\; \mathcal{M}^\circ\times\mathcal{N}^\circ \longrightarrow \Ens\;,
\tag*{(4.1.10.1)}
\]
tels que pour tout $x$ dans $\mathcal{M}$, le foncteur
\[
y \longmapsto F(x,y) \;:\; \mathcal{N}^\circ \longrightarrow \Ens
\tag*{(4.1.10.2)}
\]
soit dans $\uHom^!(\mathcal{N}^\circ,\Ens)$, i.e. transforme petites $\varinjlim$ en $\varprojlim$\,, et que, symétriquement, pour tout $y$ dans $\mathcal{N}$, le foncteur
\[
x \longmapsto F(x,y) \;:\; \mathcal{M}^\circ \longrightarrow \Ens
\tag*{(4.1.10.3)}
\]
soit dans $\uHom^!(\mathcal{M}^\circ,\Ens)$. Mais cela signifie aussi que ces foncteurs sont \textit{représentables}\,; donc on trouve aussi
\[
\mathcal{M}\boxtimes\mathcal{N} \xleftarrow{\sim} \uHom^!(\mathcal{N}^\circ,\mathcal{M})
\tag*{(4.1.11)}
\]
et
\[
\mathcal{M}\boxtimes\mathcal{N} \xleftarrow{\sim} \uHom^!(\mathcal{M}^\circ,\mathcal{N})
\tag*{(4.1.12)}
\]
\msp{page 56}
en utilisant (dans chacune de ces formules) dans deux côtés les équivalences de catégories
\[
\left\{
\begin{aligned}
&\mathcal{M} &&\xrightarrow{\sim} &&\uHom^!(\mathcal{M}^\circ,\Ens)\\
&\mathcal{N} &&\xrightarrow{\sim} &&\uHom^!(\mathcal{N}^\circ,\Ens)\;.
\end{aligned}
\right.
\tag*{(4.1.13)}
\]
On peut dire aussi que la catégorie $\mathcal{M}\boxtimes\mathcal{N}$ est équivalente à celle des \textit{couples} de foncteurs
\[
\left\{
\begin{aligned}
&\varphi &&: &&\mathcal{M}^\circ &&\longrightarrow &&\mathcal{N}\\
&\psi &&: &&\mathcal{N}^\circ &&\longrightarrow &&\mathcal{M}\;,
\end{aligned}
\right.
\tag*{(4.1.14)}
\]
munis de deux morphismes d'adjonction
\[
\left\{
\begin{aligned}
&\psi\varphi^\circ &&\xleftarrow{\alpha} &&\mathrm{id}_{\mathcal{M}}\\
&\varphi\psi^\circ &&\xleftarrow{\beta} &&\mathrm{id}_{\mathcal{N}}
\end{aligned}
\right.
\tag*{(4.1.15)}
\]
satisfaisant les deux compatibilités habituelles exprimant l'adjonction de $\varphi^\circ$ et de $\psi$ (ou, de façon équivalente, de $\psi^\circ$ et de $\varphi$). Le foncteur
\[
(\varphi,\psi) \longmapsto \varphi \qquad (\text{resp.}\;\; (\varphi,\psi) \longmapsto \psi)
\]
donne une équivalence de la catégorie de ces couples avec $\uHom^!(\mathcal{M}^\circ,\mathcal{N})$ (resp. avec $\uHom^!(\mathcal{N}^\circ,\mathcal{M})$), de sorte qu'on a un carré commutatif d'équivalences canoniques de foncteurs :

% ---- from der-ch18_p006-010.tex
\printedpage{4}
% head: A. Grothendieck : Dérivateurs, Ch. XVIII, Édition des Universités de Montpellier II et Paris VII

\msp{page 57}

\[
\begin{tikzcd}[column sep=small]
 & \cM \tbt \cN \arrow[dl, "\sim" sloped] \arrow[dr, "\sim" sloped] & \\
\uHom^{!}(\cM^{\circ},\cN) \arrow[dr, "\sim" sloped] & & \uHom^{!}(\cN^{\circ},\cM) \arrow[dl, "\sim" sloped] \\
 & \underbrace{\cM \bt \cN}_{\defeq \uHom^{!!}(\cM^{\circ}\times\cN^{\circ},\Ens)} &
\end{tikzcd}
\tag*{(4.1.16)}
\]

\texttt{[où $\cM \tbt \cN$ est la catégorie des couples définie dans la page précédente.]}

Je dis que $\cM \bt \cN$ \textit{est une catégorie pseudo-topos}. Il est clair, \texttt{[en]} outre, qu'elle admet des petites $\varprojlim$ , qui se calculent argument par argument dans $\uHom(\cM^{\circ}\times\cN^{\circ},\Ens)$, ou au choix dans $\uHom^{!}(\cM^{\circ},\cN)$, ou $\uHom^{!}(\cN^{\circ},\cM)$. Comment construire des petites $\varinjlim$ ? Il n'est pas question d'invoquer un argument de cogénérateurs. Il faut d'ailleurs trouver aussi une petite sous-catégorie strictement génératrice. Visiblement, il y a quelque chose de non trivial à démontrer !

\section*{4.2\quad Cafouillages sur $\cM \bt \cN$}

Revenons à la situation générale d'un pseudo-topos $\cX$, décrit par la catégorie pseudo-topos $\cM$. Si $\cN$ est une catégorie quelconque, on définit la catégorie

\msp{page 58}
\[
\mathrm{Fais}(\cX,\cN) \qquad \text{ou} \qquad \cF(\cX,\cN)
\]
des faisceaux sur $\cX$, à valeurs dans $\cN$, par
\[
\cF(\cX,\cN) \ \defeq\ \uHom^{!}(\cM^{\circ},\cN) \ .
\tag*{(4.2.1)}
\]
On a donc un foncteur pleinement fidèle
\[
\cF(\cX,\cN) \ \defeq\ \uHom^{!}(\cM^{\circ},\cN) \ \linj\ \Hom^{!}(\cN^{\circ},\overline{\cM}) \ ,
\tag*{(4.2.2)}
\]
où $\overline{\cM}$ désigne la sous-catégorie de $\cM^{\wedge}$ (\textbf{N.B.} $\cM^{\wedge}$ n'est pas en général une $\mathcal{U}$-catégorie qu'à cela ne tienne) formée des foncteurs \textit{représentables} $\cM^{\circ}\longrightarrow\Ens$ (laquelle est une $\mathfrak{U}$-catégorie, équivalente à $\cM$ par l'inclusion canonique $\cM \linj \overline{\cM}$). Donc on a à isomorphisme unique près une factorisation de (4.2.2) en
\[
\cF(\cX,\cN) \ \overset{\text{pl. fid.}}{\lhook\joinrel\relbar\joinrel\relbar\joinrel\longrightarrow}\ \uHom^{!}(\cN^{\circ},\cM) \ ,
\tag*{(4.2.3)}
\]
l'image essentielle étant formée des foncteurs
\[
\psi \ :\ \cN^{\circ} \longrightarrow \cM
\]
tels que pour tout $x$ dans $\cM$,
\[
y \longmapsto \Hom(y,\psi(x)) \ :\ \cN^{\circ} \longrightarrow \Ens
\]
soit représentable, i.e. des foncteurs qui
% foot: Version 17 décembre 2025
\printedpage{5}
% head: A. Grothendieck : Dérivateurs, Ch. XVIII, Édition des Universités de Montpellier II et Paris VII
\msp{page 59}
admettent un \textit{adjoint à gauche}. Quand $\cN$ lui-même est une catégorie pseudo-topos, cette condition est automatique (elle résulte de l'hypothèse que $\psi$ commute aux petites $\varprojlim$ ). Donc dans ce cas, (4.2.3) est une équivalence de catégories.

Soit maintenant $S$ ($S$ comme “site”) une \textit{petite} sous-catégorie de $\cM$, génératrice par épimorphismes stricts. On sait que cela signifie que le foncteur composé
\[
\cM \xrightarrow{\ \sim\ } \uHom^{!}(\cM^{\circ},\Ens) \xrightarrow{\text{restriction à } S} S^{\wedge} \ (\ \defeq \uHom(S^{\circ},\Ens)\ )
\tag*{(4.2.4)}
\]
est pleinement fidèle, et ceci implique plus généralement, que pour toute catégorie $\cN$, le foncteur restriction
\[
\cF(\cX,\cN) \ \defeq\ \uHom^{!}(\cM^{\circ},\cN) \ \overset{\text{pl. fid.}}{\lhook\joinrel\relbar\joinrel\relbar\joinrel\longrightarrow}\ \uHom(S^{\circ},\cN)
\tag*{(4.2.5)}
\]
est pleinement fidèle. En effet, le deuxième membre, via le plongement pleinement fidèle
\[
\cN \ \overset{\text{pl. fid.}}{\lhook\joinrel\relbar\joinrel\relbar\joinrel\longrightarrow}\ \uHom^{!}(\cN^{\circ},\Ens) \ \linj\ \cN^{\wedge} \ ,
\]
se plonge dans
\[
\uHom(S^{\circ},\uHom^{!}(\cN^{\circ},\Ens)) \ \simeq\ \Hom^{!}(\cN^{\circ},S^{\wedge}) \ ,
\]
et la pleine fidélité de (4.2.5) équivaut donc à celle du foncteur composé
\[
\cF(\cX,\cN) \ \defeq\ \uHom^{!}(\cM^{\circ},\cN) \ \longrightarrow\ \uHom^{!}(\cN^{\circ},S^{\wedge}) \ \simeq\ \uHom(S^{\circ},\uHom^{!}(\cN^{\circ},\Ens)) \ .
\tag*{(4.2.6)}
\]

\msp{page 60}
Or on a commutativité dans le carré
\[
\begin{tikzcd}[column sep=large, row sep=large]
\cF(\cX,\cN) \defeq \uHom^{!}(\cM^{\circ},\cN) \arrow[r, "\text{\scriptsize(4.2.5)}"] \arrow[d, hook, "\text{\scriptsize(4.2.3)}"', "\text{\scriptsize pleinement fidèle}"] & \uHom(S^{\circ},\cN) \arrow[d, hook, "\text{\scriptsize pleinement fidèle}"] \\
\uHom^{!}(\cN^{\circ},\cM) \arrow[r, hook, "\text{\scriptsize(4.2.4)}", "\text{\scriptsize pleinement fidèle}"'] & \underbrace{\uHom^{!}(\cN^{\circ},S^{\wedge})}_{\simeq\, \uHom(S^{\circ},\uHom^{!}(\cN^{\circ},\Ens))}
\end{tikzcd}
\tag*{(4.2.7)}
\]
(où les flèches verticales sont des équivalences, si $\cN$ est lui-même une catégorie pseudo-topos). Notons d'autre part que l'on a une équivalence
\[
\uHom(S^{\circ},\cN) \ \xleftarrow{\ \sim\ }\ \uHom^{!}(S^{\wedge\circ},\cN) \ \defeq\ \cF(\Top(S),\cN) \ ,
\tag*{(4.2.8)}
\]
où $\Top(S)$ est le pseudo-topos (en fait, un topos) défini par la catégorie pseudo-topos $S^{\wedge}$. Pour ceci, notons qu'on a
\[
\uHom^{!}(S^{\wedge\circ},\cN) \ \xrightarrow{\ \sim\ }\ \uHom_{!}(S^{\wedge},\cN^{\circ})^{\circ} \ ,
\]
et d'autre part le foncteur restriction à $S \linj S^{\wedge}$ donne
\[
\uHom_{!}(S^{\wedge},\cN^{\circ}) \ \xrightarrow{\ \sim\ }\ \uHom(S,\cN^{\circ}) \ ,
\]
d'où
% foot: Version 17 décembre 2025
\printedpage{6}
% head: A. Grothendieck : Dérivateurs, Ch. XVIII, Édition des Universités de Montpellier II et Paris VII
\msp{page 61}
\[
\uHom(S^{\wedge\circ},\cN) \ \xrightarrow{\ \sim\ }\ (\uHom(S,\cN^{\circ}))^{\circ} \ \xrightarrow{\ \sim\ }\ \uHom(S^{\circ},\cN) \ ,
\]
ce qui n'est autre que (4.2.8). Ainsi :

\textbf{Proposition 4.2.1.} \textit{Soient $\cM=\cF(\cX)$ une catégorie pseudo-topos, munie d'une sous-catégorie pleine strictement génératrice $S$, $\cN$ une $\mathfrak{U}$-catégorie quelconque, alors le foncteur “restriction à $S$” suivant est pleinement fidèle.}
\[
\cF(\cX,\cN) \ \defeq\ \uHom^{!}(\cM^{\circ},\cN) \ \longrightarrow\ \uHom(S^{\circ},\cN) \ \xleftarrow{\ \sim\ }\ \underbrace{\uHom(S^{\wedge\circ},\cN)}_{\defeq\, \cF(\Top(S),\cN)} \ ,
\tag*{(4.2.1.1)}
\]
\textit{et il s'insère dans le carré commutatif (4.2.7) de foncteurs pleinement fidèles.}

\textbf{Corollaire 4.2.2.} \textit{Supposons que $\cN=\cF(\cY)$ soit également une catégorie pseudo-topos, munie d'une sous-catégorie pleine $T$ strictement génératrice. Alors}
\[
\cF(\cX,\cN) \ \simeq\ \cF(\cY,\cM) \ \simeq\ \underbrace{\uHom^{!!}(\cM^{\circ}\times\cN^{\circ},\Ens)}_{=\,\cM\bt\cN} \ ,
\]
\textit{et le foncteur restriction}
\[
\begin{aligned}
\cM\bt\cN \ \defeq\ \uHom^{!!}(\cM^{\circ}\times\cN^{\circ},\Ens) \quad &\longrightarrow \quad \uHom(S^{\circ}\times T^{\circ},\Ens)\\
&\simeq \quad (S\times T)^{\wedge}\\
&\defeq \quad \cF(\Top(S\times T))
\end{aligned}
\tag*{(4.2.2.1)}
\]
\textit{est pleinement fidèle.}

La première assertion a été vue dans 4.1. Pour la pleine fidélité de (4.2.2.1), on note que ce foncteur s'interprète, à équivalence de catégories près, comme le composé du foncteur pleinement fidèle (4.2.1.1), et du foncteur
\msp{page 62}
similaire relatif à $\cN$ (donc lui aussi pleinement fidèle)
\[
\cF(\cY,S^{\wedge}) \ \defeq\ \uHom^{!}(\cN^{\wedge},S^{\wedge}) \ \xrightarrow{\ \text{restr. à } T\ }\ \underbrace{\uHom(T^{\circ},S^{\wedge})}_{\uHom(S^{\circ}\times T^{\circ},\Ens)} \ ,
\tag*{(4.2.2.2)} % [sic] \uHom^{!}(\cN^{\wedge},S^{\wedge}) -> \uHom^{!}(\cN^{\circ},S^{\wedge}), cf. (4.2.2.3) where the same F(Y,S^) is Hom!(N°,S^)
\]
compte tenu de l'équivalence de symétrie
\[
\cF(\cY,S^{\wedge}) \ \defeq\ \uHom^{!}(\cN^{\circ},S^{\wedge}) \ \simeq\ \uHom^{!}(S^{\wedge\circ},\cN) \ .
\tag*{(4.2.2.3)}
\]

Au total, quand $\cM$, $\cN$ sont tous deux des catégories pseudo-topos, munies de sous-catégories strictement génératrices $S$ resp. $T$, on a un plongement pleinement fidèle canonique
\[
\cM\bt\cN \ \overset{\rho}{\linj}\ (S\times T)^{\wedge} \ ,
\tag*{(4.2.9)}
\]
% foot: Version 17 décembre 2025
\printedpage{7}
% head: A. Grothendieck : Dérivateurs, Ch. XVIII, Édition des Universités de Montpellier II et Paris VII
qu'on interprète par la suite comme $(f_{S}^{\cM}\times f_{T}^{\cN})_{*}$, où
\[
\left\{
\begin{aligned}
f_{S}^{\cM} \ &:\ \cX \ \longrightarrow\ \Top S\\
f_{T}^{\cN} \ &:\ \cY \ \longrightarrow\ \Top T
\end{aligned}
\right.
\tag*{(4.2.10)}
\]
sont les morphismes de pseudo-topos, donnés par les plongements pleinement fidèles (foncteurs image directe)
\[
\left\{
\begin{aligned}
(f_{S}^{\cM})_{*} \ &:\ \cM \ \longrightarrow\ S^{\wedge}\\
(f_{T}^{\cN})_{*} \ &:\ \cN \ \longrightarrow\ T^{\wedge} \ ,
\end{aligned}
\right.
\tag*{(4.2.11)}
\]
lesquels admettent bel et bien des adjoints à gauche
\msp{page 63}
\[
\left\{
\begin{aligned}
(f_{S}^{\cM})^{*} \ &:\ S^{\wedge} \ \longrightarrow\ \cM\\
(f_{T}^{\cN})^{*} \ &:\ T^{\wedge} \ \longrightarrow\ \cN \ ,
\end{aligned}
\right.
\tag*{(4.2.12)}
\]
à savoir les prolongements canoniques (par commutativité aux petites $\varinjlim$ ) des foncteurs d'inclusion $S \linj \cM$, $T \linj \cN$.

La question essentielle, à présent, c'est visiblement de construire un adjoint à gauche pour le foncteur (4.2.9). Il revient au même, pour $(s,t)$ dans $S\times T$, de définir un objet
\[
s\bt t \ \in\ \Ob\cM\bt\cN \ ,
\tag*{(4.2.13)}
\]
tel que l'on ait, pour $F$ variable dans $\cM\bt\cN$, un isomorphisme fonctoriel en $F$,
\[
\underbrace{\rho(F)(s,t)}_{\defeq\, F(s,t)} \ \simeq\ \Hom(s\bt t,F) \ .
\]
Donc il faut prouver ceci :

\textbf{Lemme 4.2.3} (\footnote[4]{Pas prouvé ; peut-être faux !}). \textit{Soient $\cM$, $\cN$ deux catégories pseudo-topos, et $s\in\Ob\cM$, $t\in\Ob\cN$, et considérons sur $\cM\bt\cN\defeq\uHom^{!!}(\cM^{\circ}\times\cN^{\circ},\Ens)$ le foncteur $F\longmapsto F(s,t)$. Ce foncteur est} représentable \textit{(par un objet que je note $s\bt t$).}

L'idée naturelle serait de définir $s\bt t$ par
\[
(s\bt t)(x,y) \ =\ s(x)\times t(y) \ =\ \Hom_{\cM}(x,s)\times\Hom_{\cN}(y,t) \ ,
\]

\msp{page 64}
ce qui définit bien un foncteur
\[
\cM^{\circ}\times\cN^{\circ} \longrightarrow \Ens \ .
\]
Malheureusement, pour un des arguments fixé, disons pour $y$ fixé, ce foncteur n'est \textit{pas} compatible aux limites projectives, p.ex. il ne transforme pas l'objet final de $\cM^{\circ}$ correspondant à l'objet initial de $\cM$) en objet final de Ens, mais en $t(y)=\Hom_{\cN}(y,t)$. % [sic] "de $\cM$)" -> closing parenthesis without opening one, unbalanced as printed
% foot: Version 17 décembre 2025
\printedpage{8}
% head: A. Grothendieck : Dérivateurs, Ch. XVIII, Édition des Universités de Montpellier II et Paris VII
Néanmoins, je pense que $s\bt t$ doit bel et bien se définir comme produit de deux objets de $\cM\bt\cN$, dépendant fonctoriellement de $s$ resp. de $t$ (\footnote[5]{également faux, cf. plus bas.}). On est donc amené dès à présent, qu'on le veuille ou non, à essayer de définir des foncteurs canoniques, commutant aux petites $\varinjlim$ ,
\[
\cM \longrightarrow \cM\bt\cN \longleftarrow \cN
\tag*{(4.2.13)} % [sic] (4.2.13) -> (4.2.14), the label (4.2.13) is already used for s⊠t ∈ Ob M⊠N
\]
(lesquels seront les foncteurs image inverse pour les projections canoniques
\begin{tikzcd}[row sep=small, column sep=small, baseline=(\tikzcdmatrixname-2-1.base)]
 & \cX \\
\cX\times\cY \arrow[ur] \arrow[dr] & \\
 & \cY
\end{tikzcd}
) (\footnote[6]{Le cas où $\cM$, $\cN$ sont des catégories-topos, nous montre que ces foncteurs devraient être décrits comme $x\longmapsto x\bt e_{\cN}$, $y\longmapsto e_{\cM}\bt y$. Mais il n'est pas clair qu'on ait en dehors du cas des topos, et même en admettant l'existence du $x\bt y$, la formule (valable pour un topos) $x\bt y\simeq(x\bt e_{\cN})\times(e_{\cM}\bt y)$.}).

\section*{4.3\quad L'opération $\cM\bt\cN$ et l'axiome d'accessibilité}

J'ai passé la journée hier, essayant d'établir l'existence des $x\bt y$ dans $\cM\bt\cN$, sans y parvenir – du moins pas sans faire une hypothèse supplémentaire d'accessibilité,
\msp{page 65}
dont il devient évident à présent qu'il faut la faire de toutes façons. Je vais quand même énoncer le positif de mes perplexités.

\textbf{Proposition-perplexité 4.3.1.} \textit{Soient $\cM$, $\cN$ deux catégories pseudo-topos, d'où $\cM\bt\cN$ comme dans 4.1. Soient $S\subset\cM$, $T\subset\cN$ deux petites sous-catégories strictement génératrices, et considérons l'inclusion pleinement fidèle correspondante :}
\[
\begin{tikzcd}[column sep=large]
\underbrace{\cM\bt\cN}_{\defeq\, \uHom^{!!}(\cM^{\circ}\times\cN^{\circ},\Ens)} \arrow[r, hook, "\alpha"'] & S^{\wedge}\bt T^{\wedge} \ \simeq\ (S\times T)^{\wedge}. \arrow[l, dotted, bend right, "\varphi\,?"']
\end{tikzcd}
\tag*{(4.3.1.1)}
\]
\textit{Considérons les conditions suivantes :}
\begin{enumerate}[label={}]
\item[\textit{a)}] \textit{$\cM\bt\cN$ est un pseudo-topos.}
\item[\textit{b)}] \textit{Pour $x\in\Ob\cM$, $y\in\Ob\cN$, $x\bt y\in\Ob\cM\bt\cN$ existe, i.e. le foncteur (commutant aux petites $\varprojlim$ )}
\[
F \ \longmapsto\ F(x,y) \ , \qquad \cM\bt\cN \longrightarrow \Ens
\tag*{(4.3.1.2)}
\]
\textit{est représentable.}
\item[\textit{b$'$)}] \textit{Comme b), avec $x\in\Ob S$, $y\in\Ob T$.}
\item[\textit{c)}] \textit{Le foncteur $\alpha$ (lequel commute aux petites $\varprojlim$ ) admet un adjoint à gauche}
\[
\varphi \ :\ (S\times T)^{\wedge} \longrightarrow \cM\bt\cN \ .
\]
\end{enumerate}
\textit{On a alors les implications}
\[
b \ \Longrightarrow\ b' \ \Longleftrightarrow\ c \ \Longrightarrow\ a \ .
\tag*{(4.3.1.3)}
\]
% foot: Version 17 décembre 2025

% ---- from der-ch18_p011-015.tex
\setcounter{footnote}{6}

\printedpage{9}
% head: A. Grothendieck : Dérivateurs, Ch. XVIII, Édition des Universités de Montpellier II et Paris VII

\noindent\textit{De plus, si c est satisfait, on a ceci : pour $s\in S$, $t\in T$, considérons}
\[
(s,t)\ \in\ (S\times T)\ \subset\ (S\times T)^\wedge\ \simeq\ S^\wedge\boxtimes T^\wedge\ ,
\]
\textit{qui s'explicite aussi comme}
\[
\left\{
\begin{aligned}
&\rlap{$\displaystyle (s,t)\ =\ \pr_1^*(s)\times\pr_2^*(t)\ \simeq\ s\overset{S^\wedge,T^\wedge}{\boxtimes}t\in\Ob\,(S\times T)^\wedge\ ,$}\\
&\textit{i.e.}\ (s,t)(s',t')\ =\ s(s')\times t(t') &&=\ \Hom_S(s',s)\times\Hom_T(t',t)\\
& &&=\ \Hom_{S\times T}((s',t'),(s,t))\ ,
\end{aligned}
\right.
\tag*{(4.3.1.4)}
\]
\[
\begin{tikzcd}
S\times T \arrow[r, "\mathrm{pr}_2"] \arrow[d, "\mathrm{pr}_1"] & T\\
S\ ,
\end{tikzcd}
\]
\msp{page 66}\textit{on a}
% note: the label (4.3.1.4) is printed for this display as well as for the preceding one
\[
s\overset{\mathcal M,\mathcal N}{\boxtimes}t\ \simeq\ \varphi(s,t)\ ,
\tag*{(4.3.1.4)}
\]
\textit{et $\varphi(S\times T)=\{s\boxtimes t\mid s\in S,\ t\in T\}$ est une partie strictement génératrice de $\mathcal M\boxtimes\mathcal N$ (}\footnote{\textbf{N.B.} On a bien sûr $(s,t)=s\overset{S^\wedge,T^\wedge}{\boxtimes}t$.}\textit{).}

\noindent\textbf{N.B.} Je ne connais pas d'exemple où les conditions envisagées ne soient satisfaites, mais doute pourtant à présent que même la plus faible le soit toujours, i.e. que $\mathcal M\boxtimes\mathcal N$ soit un pseudo-topos dès que $\mathcal M$ et $\mathcal N$ le sont.

\noindent\textsc{Démonstration de} 4.3.1. L'implication b $\Longrightarrow$ b$'$ est tautologique. L'implication b$'$ $\Longleftrightarrow$ c résulte du

% marginal vertical bar printed beside Lemme 4.3.2 (from "Lemme 4.3.2." to "soit représentable.")
\noindent\textbf{Lemme 4.3.2.} \textit{Soient $\mathcal E$, $X$ des $\mathfrak U$-catégories avec $X$ petite, $\mathcal E$ stable par petites $\varprojlim$ ,}
\[
\alpha\ :\ \mathcal E\longrightarrow X^\wedge
\]
\textit{un foncteur. Pour que $\alpha$ admette un adjoint à gauche $\varphi$, il faut et il suffit que pour tout $x\in\Ob X$, le foncteur}
\[
\left\{
\begin{array}{ccl}
\xi & \longmapsto & \Hom_{X^\wedge}(x,\alpha(\xi))\ =\ \alpha(\xi)(x)\\
\mathcal E^\circ & \longrightarrow & \mathrm{Ens}
\end{array}
\right.
\]
\textit{soit représentable.}

\noindent C'est évidemment nécessaire. La suffisance résulte du fait que $X\subset X^\wedge$ est stricte%
% foot: Version 17 décembre 2025
\printedpage{10}%
% head: A. Grothendieck : Dérivateurs, Ch. XVIII, Édition des Universités de Montpellier II et Paris VII
\msp{page 67}%
ment génératrice dans $X^\wedge$, i.e. que tout objet $F$ de $X^\wedge$ est limite inductive $F=\varinjlim\limits_I F_i$, les $F_i$ dans $X$ (on prend $I=X/F$), d'où
\[
\begin{aligned}
\Hom(F,\alpha(\xi))\ &\simeq\ \varprojlim_I\ \underbrace{\Hom(F_i,\alpha(\xi))}_{\Hom(\rho(F_i),\xi)}\\
&\simeq\ \Hom(\varinjlim_I\,(\rho(F_i),\xi))\ ,
\end{aligned}
\]
donc le foncteur $\xi\longmapsto\Hom(F,\alpha(\xi))$ est représenté par $\rho(F)\overset{\text{déf}}{=}\varinjlim\limits_I\rho(F_i)$.

\noindent Prouvons c $\Longrightarrow$ a. Pour l'existence des petites $\varinjlim$ dans $\mathcal M\boxtimes\mathcal N$, et l'existence d'une partie strictement génératrice, on invoque le lemme général :

% marginal vertical bar printed beside Lemme 4.3.3 (from "Lemme 4.3.3." to "comme")
\noindent\textbf{Lemme 4.3.3.} \textit{Soit $\alpha:\mathcal E\hookrightarrow\mathcal F$ un foncteur pleinement fidèle admettant un adjoint à gauche $\rho$.}
\begin{enumerate}[label={}]
\item[\textcircled{a}] \textit{Si les $\varinjlim$ d'un type donné $I$ existent dans $\mathcal F$, elles existent dans $\mathcal E$, et se calculent comme}
\[
\varinjlim_I{}_{\mathcal E}\,\xi_i\ =\ \rho(\varinjlim_I{}_{\mathcal F}\,\alpha(\xi_i))\ .
\tag*{(4.3.3.1)}
\]
\item[\textcircled{b}] \textit{Si $U$ est une petite partie strictement génératrice de $\mathcal F$, $\rho(U)$ est une petite partie strictement génératrice dans $\mathcal F$.}
\end{enumerate}

\noindent En effet, si $\xi$ désigne le deuxième membre \texttt{[de l'égalité 4.3.3.1]}, et $\eta$ un objet variable de $\mathcal E$, on aura
\[
\begin{aligned}
\Hom_{\mathcal E}(\xi,\eta)\ &\overset{\text{adj.}}{\simeq}\ \Hom_{\mathcal F}(\varinjlim_I\alpha(\xi_i),\alpha(\eta))\\
&\simeq\ \varprojlim_I\quad\underbrace{\Hom_{\mathcal F}(\alpha(\xi_i),\alpha(\eta))}_{\simeq\,\Hom_{\mathcal E}(\xi_i,\eta)\ \ \text{car $\alpha$ pleinement fidèle}}\quad .
\end{aligned}
\]
\msp{page 68}
\[
\simeq\ \ \varprojlim_I\Hom_{\mathcal E}(\xi_i,\eta)\ ,
\]
ce qui prouve que $\xi$ est une limite inductive des $\xi_i$ dans $\mathcal E$ (\footnote{Pour prouver b), on note que la pleine fidélité de $\alpha$ implique que $\rho$ est essentiellement surjectif ; comme $\rho$ commute aux petites $\varinjlim$ et que tout objet de $\mathcal F$ est de la forme $\varprojlim\limits_I x_i$, les $x_i$ dans $U$, on aura $\rho(x)=\varinjlim\limits_I\rho(x_i)$, les $\rho(x_i)$ dans $\rho(U)$, on gagne.}).

\noindent Ceci prouve que (moyennant c)) $\mathcal M\boxtimes\mathcal N$ est stable par petites $\varinjlim$ , donc il reste à prouver que $\mathcal M\boxtimes\mathcal N$ admet une petite famille strictement génératrice. Cela résultera du b) dans le lemme, puisque $S\times T$ engendre strictement $(S\times T)^\wedge$. Cela prouve en même temps la dernière assertion de la proposition, sous réserve de prouver la formule (4.1.3.4). Il faut donc prouver un isomorphisme fonctoriel en $F\in\Ob\mathcal M\boxtimes\mathcal N$, % [sic] (4.1.3.4) -> (4.3.1.4), the formula s ⊠ t ≃ φ(s,t) to be proved is labelled (4.3.1.4) on p. 9
\[
\Hom(\rho(s,t),F)\ \simeq\ F(s,t)\ ,
\]
mais par adjonction on a
\[
\Hom(\rho(s,t),F)\ \simeq\ \Hom((s,t),\alpha(F))\ =\ \alpha(F)(s,t)\ =\ F(s,t)\ ,
\]
OK.

% foot: Version 17 décembre 2025
\printedpage{11}
% head: A. Grothendieck : Dérivateurs, Ch. XVIII, Édition des Universités de Montpellier II et Paris VII

\noindent\msp{page 69}\textbf{Proposition 4.3.4.} \textit{Soit $\mathcal M$ une $\mathfrak U$-catégorie. Les conditions suivantes sont équivalentes.}
\begin{enumerate}[label={}]
\item[\textit{a)}] \textit{$\mathcal M$ est une catégorie $\mathfrak U$-pseudo-topos, i.e. elle est stable par $\mathfrak U$-limites inductives, et admet une $\mathfrak U$-petite sous-catégorie $S$ strictement génératrice.}
\item[\textit{b)}] \textit{Il existe une $\mathfrak U$-petite catégorie $S$, et un foncteur $i:\mathcal M\hookrightarrow S^\wedge$} pleinement fidèle, \textit{ayant un adjoint à gauche.}
\end{enumerate}

\noindent\textsc{Démonstration}. On a a) $\Longrightarrow$ b), car il suffit de prendre pour $S$ dans b) celui spécifié dans a), l'hypothèse sur $S$ implique que le foncteur canonique
\[
\mathcal M\ \longrightarrow\ S^\wedge\ ,\qquad\quad \text{composé de } \mathcal M\hookrightarrow\ \mathcal M^\wedge\ \xrightarrow{\ \text{restriction}\ }S^\wedge
\]
est pleinement fidèle. Qu'il admette un adjoint à gauche résulte du lemme 4.3.2. L'implication b) $\Longrightarrow$ a) résulte du lemme 4.3.3.

\noindent\textbf{Corollaire 4.3.5.} \textit{Soit $\mathcal M$ une $\mathfrak U$-catégorie. Pour que $\mathcal M$ soit un $\mathfrak U$-pseudo-topos} accessible \textit{(i.e. dont tous les objets sont accessibles), il faut et il suffit qu'il existe $S$ dans} Cat \textit{et un foncteur pleinement fidèle et} accessible $\mathcal M\overset{\alpha}{\hookrightarrow}S^\wedge$\textit{, admettant un adjoint à gauche $\rho$.}

\noindent\msp{page 70}\textsc{Démonstration}. Il suffit de noter que si $S$, $\alpha$ sont comme dans 4.3.4 b), alors $\alpha$ est accessible si et seulement si les $\rho(S)$ sont formées d'objets accessibles dans $\mathcal M$ ; comme cette partie est strictement génératrice, on sait que cela implique que \textit{tout} objet de $\mathcal M$ est accessible. (Cf. SGA 4 I 9.9.)

\noindent Je vais essayer de donner un énoncé plus complet :

\noindent\textbf{Théorème 4.3.6.} \textit{Soit $\mathcal M$ une $\mathfrak U$-catégorie. Conditions équivalentes :}
\begin{enumerate}[label={}]
\item[\textit{a)}] \textit{$\mathcal M$ stable par petites $\varinjlim$ , et elle admet une petite sous-catégorie pleine $S$, strictement génératrice, et formée d'objets accessibles de $\mathcal M$.}
\item[\textit{a$'$)}] \textit{$\mathcal M$ est une catégorie $\mathfrak U$-pseudo-topos faiblement accessible, i.e. satisfait a), avec en plus} tout \textit{élément de $\mathcal M$ (pas seulement ceux de $S$) accessible dans $\mathcal M$.}
\item[\textit{b)}] \textit{Il existe une petite catégorie $S$, et un foncteur $\alpha:\mathcal M\hookrightarrow S^\wedge$} pleinement fidèle et accessible, \textit{tel que $\alpha$ admette un adjoint à gauche.}
\item[\textit{c)}] \msp{page 71}\textit{$\mathcal M$ est stable par petites $\varinjlim$ , est accessible, et admet une filtration cardinale} $(\Filt^\pi\mathcal M)_{\pi\geqslant\pi_0}$.
\item[\textit{d)}] \textit{Il existe un cardinal $\pi_1$, une petite catégorie $\mathcal M_{\pi_1}$, stable par limites inductives de cardinal $\leqslant\pi_1$, tels qu'on ait une équivalence de catégories}
\[
\mathcal M\ \xleftarrow{\ \sim\ }\ \Ind_{\pi_1}\mathcal M_{\pi_1}\ ,
\tag*{(4.3.6.1)}
\]
\textit{où le deuxième membre désigne la sous-catégorie pleine de} $\Ind(\mathcal M_{\pi_1})$ \textit{formée des ind-objets $(X_i)_{i\in I}$ de $\mathcal M_{\pi_1}$ tels que $I$ soit grand devant $\pi_1$.}
\end{enumerate}

\noindent\textsc{Démonstration de} 4.3.6 (\footnote{Pour une démonstration complète (celle donnée ici canule à partir d'un certain moment), cf. §4.6 (p. 139-148).}). On a déjà vu a$'$ $\Longrightarrow$ a $\Longrightarrow$ b $\Longrightarrow$ a$'$, i.e. l'équivalence de a, a$'$, b, soit PTA cette condition. Nous allons prouver les implications
\[
\mathrm{PTA}\ \Longrightarrow\ \mathrm{c}\ \Longrightarrow\ \mathrm{d}\ \Longrightarrow\ \mathrm{PTA}\ ,
\]
% foot: Version 17 décembre 2025
\printedpage{12}%
% head: A. Grothendieck : Dérivateurs, Ch. XVIII, Édition des Universités de Montpellier II et Paris VII
ce qui établira le théorème.

\noindent\textit{Prouvons PTA $\Longrightarrow$ c.} Soit $S$ une petite sous-catégorie pleine strictement génératrice. Ainsi, \msp{page 72}pour tout $x$ dans $\mathcal M$ on a
\[
x\ =\ \varinjlim_{s\in S/x}s
\tag*{(4.3.6.2)}
\]
(cf. SGA 4 I 7.2 (i)). Soit $\pi_0=\card\Fl(S)$, et pour tout cardinal $\pi\geqslant\pi_0$, soit
\[
\begin{array}{lcl}
\Filt^\pi\mathcal M & = & \text{sous-catégorie pleine de $\mathcal M$ formée des objets $x$ tels que}\\
 & & \card\,(\Fl(S/x))\leqslant\pi\ .
\end{array}
\tag*{(4.3.6.3)}
\]
Je veux prouver que l'on trouve ainsi une filtration cardinale sur $\mathcal M$, i.e. que les conditions a) b) c) de SGA 4 I 9.12 sont satisfaites. C'est prouvé dans loc. cit. 9.13 sous des hypothèses légèrement différentes, la condition d'accessibilité de $\mathcal M$ étant remplacée par celle (beaucoup moins sympathique!) que les familles de flèches épimorphiques strictes dans $\mathcal M$ sont épimorphiques strictes universelles. Cette condition n'est utilisée que dans la démonstration de la condition b), à savoir la stabilité de $\Filt^\pi\mathcal M$ par $\varinjlim\limits_I$ de cardinal $\leqslant\pi$ (sans avoir d'ailleurs à supposer $I$ filtrante, ni a fortiori grand devant $\pi_0$). Il faut d'ailleurs dire que la définition de $\Filt^\pi\mathcal M$ donnée dans loc. cit. n'est pas tout à fait la même que (4.3.6.3), \msp{page 73}mais c'est
% the printed line break "car-/dinal" inside the display is kept as printed
\[
\begin{array}{lcl}
\Filt'^{\pi}\mathcal M & = & \text{sous-catégorie pleine de $\mathcal M$ formée des objets $x$ tels qu'il}\\
 & & \text{existe une famille épimorphique stricte $x_i\xrightarrow{\alpha_i}x$ de car-}\\
 & & \text{dinal $\leqslant\pi$, avec $x_i\in S\ \forall\ i\in I$.}
\end{array}
\tag*{(4.3.6.3$'$)}
\]
Ainsi on a
\[
\Filt^\pi\mathcal M\ \subset\ \Filt'^{\pi}\mathcal M\ ,
\]
mais dans loc. cit. on prouve que pour $x$ dans $\Filt'^{\pi}\mathcal M$, on a
\[
\card\Fl(S/x)\ \leqslant\ \pi^{\pi_0}\ ,
\]
d'où
\[
\Filt^\pi\mathcal M\ \subset\ \Filt'^{\pi}\mathcal M\ \subset\ \Filt^{\pi^{\pi_0}}\mathcal M\ ,
\]
donc si $\pi=\pi^{\pi_0}$, p.ex. $\pi=2^c$ avec $c\geqslant\pi_0$, alors on trouve
\[
\Filt^\pi\mathcal M\ =\ \Filt'^{\pi}\mathcal M\ .
\]

\noindent Donc les deux filtrations on l'air $\pm$ équivalentes. Mais la démonstration donnée de la stabilité de $\Filt'^{\pi}$ semble vraiment faire appel de façon essentielle à la condition d'universalité des familles épimorphiques strictes (donc épimorphiques effectives, avec les hypothèses faites impliquant l'existence de petites $\varprojlim$ ). Donc il me faut bon gré, mal gré, refaire le travail! % [sic] on l'air -> ont l'air, verb avoir (3rd person plural) intended

% foot: Version 17 décembre 2025
\printedpage{13}
% head: A. Grothendieck : Dérivateurs, Ch. XVIII, Édition des Universités de Montpellier II et Paris VII

\noindent\msp{page 74}Condition a) : $\Filt^\pi$ équivalente à une petite catégorie. C'est immédiat, puisque par (4.3.6.2), $\Filt^\pi$ est formée de $\varprojlim\limits_I$ dans $\mathcal M$ d'objets de $S$ (petite catégorie), avec $I$ telle que $\card\Fl I\leqslant\pi$.

\noindent Condition b) : $\Filt^\pi$ stable par $\varprojlim\limits_I$ , quand $\card\Fl I\leqslant\pi$ (\footnote{quand $I$ filtrante et grande devant $\pi_0$ -- je ne le prouve que dans ce cas, hélas!}). Soit donc $(x_i)_{i\in I}$, où
\[
\begin{array}{ccc}
I & \longrightarrow & \mathcal M\\[1ex]
i & \longmapsto & x_i
\end{array}
\]
un système inductif dans $\mathcal M$, avec des $x_i$ dans $\Filt^\pi$. Considérons le produit fibré
\[
\begin{tikzcd}
I \arrow[d, "\varphi"'] & S/I = J \arrow[l, "\text{cofibré}"'] \arrow[d] \arrow[dr, dotted, "\psi"] & & \\
\mathcal M & S/\mathcal M \arrow[l, "\text{but}"', "\text{Cat-cofibré}"] \arrow[r, "\text{source}"] & S \arrow[r, hook] & \mathcal M\ ,
\end{tikzcd}
\tag*{(4.3.6.4)}
\]
où $J$ est Cat-cofibré sur $I$, sa fibre en $i\in\Ob I$ étant isomorphe à $S/x_i$, donc $\card\Fl J_i\leqslant\pi$. Comme par hypothèse $\card\Fl I\leqslant\pi$, on en conclut que $\card\Fl J\leqslant\pi^3=\pi$. (\textbf{N.B.} L'ensemble des flèches $v:j_0\longrightarrow j_1$ de $J$ au dessus d'une flèche $u:i_0\longrightarrow i_1$ de $I$ est “contenu” dans l'ensemble produit $\Ob J_{i_0}\times\Fl J_{i_1}$, par l'application
\[
v\longmapsto(\operatorname{source}v,v':u_*(j_0)\longrightarrow j_1)\ ,
\]
donc de cardinal majoré par $\pi^2$.) On \msp{page 75}a alors
\[
x\ \overset{\text{déf}}{=}\ \varinjlim_I x_i\ =\ \varinjlim_J\ \underbrace{\xi_j}_{=\,\psi(j)}\qquad\quad\mbox{(\footnotemark)}\ ,
\]
\footnotetext{On utilise la propriété de distributivité des $\varinjlim$ , relative à un foncteur Cat-cofibrant
\[
\begin{tikzcd}[ampersand replacement=\&]
J \arrow[r] \arrow[dr] \& I\\
\& \mathcal M\ ,
\end{tikzcd}
\]
explicitée dessus.}%
où cette fois la limite est une limite d'objets de $S$, prise dans $\mathcal M$. Il s'ensuit que $x\in\Filt'^{\pi}$, ce n'est pas tout à fait $x\in\Filt^\pi$. (Ça implique $x\in\Filt^{\pi^{\pi_0}}$, donc $x\in\Filt^\pi$ si $\pi=\pi^{\pi_0}$, p.ex. $\pi=2^c$ avec $c\geqslant\pi_0$.) Mais c'est le moment d'utiliser le fait que les objets de $S$ sont accessibles, et on prendra $\pi_0$ tel que les objets de $S$ soient $\pi_0$-accessibles, et cette fois on prendra quand même $I$ filtrante et grande devant $\pi_0$ (pour nous borner à ce qui est exigé dans la condition b) de loc. cit.). On doit prouver que $\card\,(\Fl(S/x))\leqslant\pi$ (ce qui signifie $x\in\Filt^\pi$). Mais comme les objets de $S$ sont $\pi_0$-accessibles et $I$ grande devant $\pi_0$, on voit aussitôt que
\[
S/x\ \simeq\ \varinjlim_i S/x_i\ ,
\]
% foot: Version 17 décembre 2025

% ---- from der-ch18_p016-020.tex
\setcounter{footnote}{11}

\printedpage{14}
% head: A. Grothendieck : Dérivateurs, Ch. XVIII, Édition des Universités de Montpellier II et Paris VII

d'où
\[
\Fl\,S/x \;\simeq\; \varinjlim_{i} \Fl\,S/x_i \;,
\]
donc
\[
\card\Fl(S/x) \;\leqslant\; (\underbrace{\card I}_{\leqslant\pi})\underbrace{(\sup_{i}\card\Fl(S/x_i))}_{\leqslant\pi} \;\leqslant \pi^2 = \pi \;,
\]

\hfill q.e.d.

\msp{page 76}
% margin: hand-drawn bracket mark in the left margin beside the next two lines
Prouvons c) (\footnote{\textit{Démonstration canulée}, ne marche que si $\pi^{\pi_0} = \pi$\,!!}), i.e. que pour tout $\pi' > \pi \geqslant \pi_0$ et tout $x \in \Ob\Filt^{\pi'}$, on peut écrire
\[
x \;=\; \varinjlim_{I} x_i \qquad I \text{ grande devant } \pi\text{, les } x_i \text{ dans } \Filt^{\pi}\text{, } \card I \leqslant \pi'^{\pi}. \tag*{(4.3.6.5)}
\]
% margin: hand-drawn bracket mark in the left margin beside the next three lines
La démonstration se fait comme dans loc. cit. (p. 151-152), en changeant légèrement la définition de $I$, comme désignant l'ensemble des sous-catégories pleines $i$ de $S/x$ qui sont filtrantes, de cardinal $\leqslant \pi$ \textit{et de plus grandes devant} $\pi_0$ (\footnote{Il n'est pas sûr que $I$ soit grande devant $\pi$, sauf si on suppose $\pi^{\pi_0} = \pi$\,!}). (On pourra supposer $S/x$ filtrante en augmentant $S$ de façon à ce que $S$ soit stable dans $\mathcal{M}$ par petites $\varinjlim$ finies, donc les catégories $S/x$ aussi, ce qui implique qu'elles sont filtrantes. Il est \texttt{[phrase incomplète]})

On a donc prouvé PTA $\Longrightarrow$ c).

\textit{Prouvons c) $\Longrightarrow$ d).} Soit
\[
\pi_1 \;=\; 2^c \qquad \text{avec} \qquad c \;\geqslant\; \pi_0 \;.
\]
On sait alors que
\[
\Filt^{\pi_1}\mathcal{M} \;=\; \mathcal{M}_{\pi_1} \;\overset{\text{déf}}{=}\; \text{sous-catégorie pleine de $\mathcal{M}$ des objets $\pi_1$-accessibles de $\mathcal{M}$}
\]
(loc. cit. 9.16), ce qui implique notamment que $\mathcal{M}_{\pi_1}$ est stable par $\varinjlim$ de cardinal $\leqslant \pi_1$, en vertu de loc. cit. 9.9 (page 145). De plus,
\msp{page 77}
par loc. cit. 9.18 b) (p. 156), on a
\[
\Ind_{\pi_1}(\mathcal{M}_{\pi_1}) \xrightarrow{\;\sim\;} \mathcal{M} \;,
\]
ce qui prouve d).

\textit{Prouvons d) $\Longrightarrow$ PTA.} Par loc. cit. 9.18 a) (p. 156), on voit que les objets de $\mathcal{M}_{\pi_1}$ sont $\pi_1$-accessibles dans $\Ind_{\pi_1}(\mathcal{M}_{\pi_1})$, d'autre part il est clair que $\mathcal{M}_{\pi_1}$ est une sous-catégorie pleine strictement génératrice de $\Ind_{\pi_1}(\mathcal{M}_{\pi_1})$, donc il reste à voir que $\Ind_{\pi_1}(\mathcal{M}_{\pi_1})$ est stable par petites $\varinjlim$ , i.e. le

\textbf{Lemme 4.3.7.} \textit{Soit $\mathcal{M}_\pi$ une $\mathfrak{U}$-catégorie stable par limites inductives de cardinal $\leqslant \pi$, où $\pi$ est un cardinal infini donné. Alors} $\Ind_\pi\mathcal{M}_\pi$ \textit{est stable par petites limites inductives. De plus, $\mathcal{M} = \Ind_\pi\mathcal{M}_\pi$, munie de $\mathcal{M}_\pi \longrightarrow \mathcal{M}$, est solution d'un problème $2$-universel, à savoir
\printedpage{15}
% head: A. Grothendieck : Dérivateurs, Ch. XVIII, Édition des Universités de Montpellier II et Paris VII
que pour toute catégorie $\mathcal{N}$ stable par petites limites inductives, le foncteur restriction induit une équivalence}
\[
\uHom_{!}(\mathcal{M},\mathcal{N}) \xrightarrow{\;\sim\;} \uHom_{!\,\pi}(\mathcal{M}_\pi,\mathcal{N}) \;, \tag*{(4.3.7.1)}
\]
\textit{où le signe ! (resp. $!\,\pi$) désigne la sous-catégorie pleine de la catégorie} $\uHom$\textit{, formée des foncteurs qui commutent aux petites} $\varinjlim$ \textit{(resp. aux limites inductives de cardinal $\leqslant \pi$).}

\msp{page 78}
Prouvons d'abord que les $\varinjlim$ de cardinal $\leqslant \pi$ \textit{dans} $\mathcal{M}_\pi$ sont aussi des $\varinjlim$ \textit{dans} $\mathcal{M} = \Ind_\pi\mathcal{M}_\pi$ (\footnote{Mais \textit{attention}, $\mathcal{M}_\pi \longrightarrow \Ind(\mathcal{M}_\pi)$ ne commute \textit{pas} à ces $\varinjlim$ !}). On a donc
\[
x \xleftarrow{\;\sim\;} \varinjlim_{I} x_i \qquad \textit{dans } \mathcal{M}_\pi \;,
\]
i.e.
\[
\Hom(x,y) \xrightarrow{\;\sim\;} \varprojlim_{I} \Hom(x_i,y) \qquad \text{si } y \in \Ob\mathcal{M}_\pi \;, \tag*{$(*)$}
\]
prouvons que cela reste vrai pour $y$ dans $\mathcal{M}$, i.e. $y$ de la forme
\[
y \;=\; \varinjlim_{J} y_j \;, \qquad J \text{ grande devant } \pi\text{, les } y_j \text{ dans } \mathcal{M}_\pi.
\]
On a en effet
\[
\begin{aligned}
\Hom(x,y) \quad &\simeq \quad \varinjlim_{J} \Hom(x,y_j)\\
&\overset{(*)}{\simeq} \quad \varinjlim_{J}\, \varprojlim_{I} \Hom(x_i,y_j)\\
&\overset{\text{SGA 4 I 9.8}}{\simeq} \quad \varprojlim_{I}\, \varinjlim_{J} \Hom(x_i,y_j)\\
&\simeq \quad \varprojlim_{I} \Hom(x_i,y) \;,
\end{aligned}
\]
au total
\[
\Hom(x,y) \;\simeq\; \varprojlim_{I} \Hom(x_i,y),
\]
q.e.d.

Pour prouver que les petites $\varinjlim$ sont représentables dans $\mathcal{M} = \Ind_\pi(\mathcal{M}_\pi)$, comme elles le sont dans $\Ind(\mathcal{M}_\pi)$ (SGA 4 I 8.9.5 b, on utilise l'hypothèse que les $\varinjlim$ \textit{finies} sont représentables dans $\mathcal{M}_\pi$), il reste à voir que $\Ind_\pi(\mathcal{M}_\pi)$ est stable dans $\mathcal{M}_\pi$ par ces limites.

\msp{page 79}

\section*{4.4\quad Étude de $\Ind_\pi(C)$ (pour $C$ petite catégorie stable par $\varinjlim$ de cardinal $\leqslant \pi$)}

\textbf{4.4.1.} Finalement, la démonstration du lemme 4.3.7 n'est pas tellement évidente que je l'anticipais, et il me faut lui consacrer des préliminaires techniques, que je vais réunir ici. Dans toute la suite, $C$ représente une $\mathfrak{U}$-catégorie équivalente à une petite catégorie.

\printedpage{16}
% head: A. Grothendieck : Dérivateurs, Ch. XVIII, Édition des Universités de Montpellier II et Paris VII

\textbf{Proposition 4.4.2.} \textit{Soit $F \in C^{\wedge}$ un foncteur ind-représentable, i.e. tel que $C/F$ soit filtrante. Soit d'autre part $i \longmapsto x_i$ un foncteur $I \longrightarrow C$, où $I$ est une petite catégorie filtrante, et soit}
\[
\alpha \;:\; x \;=\; \varinjlim_{I}{}_{C^{\wedge}} x_i \longrightarrow F \qquad \text{(\footnotemark)} \tag*{(4.4.2.1)}
\]
\footnotetext{La donnée de $\alpha$ équivaut aussi à celle d'un homomorphisme d'ind-objets $(x_i)_{i\in I} \longrightarrow F$, $F$ étant considéré comme ind-objet (généralisé) indexé par $C/F$.}
\textit{un système projectif de flèches $x_i \longrightarrow F$ (}\textbf{N.B.} \textit{la limite inductive $x$ est prise dans $C^{\wedge}$, et est calculée argument par argument), ce qui équivaut à une factorisation $\varphi_\alpha : I \longrightarrow C/F$ de $I \longrightarrow C$ en}
\[
I \xrightarrow{\;\varphi_\alpha\;} C/F \hookxrightarrow[\psi]{\textit{can.}} C \;.
\]
\textit{Ceci dit, pour que $\alpha$ (4.4.2.1) soit un isomorphisme dans} $\Ind'(C) \subset C^{\wedge}$ \textit{(sous-catégorie pleine des foncteurs ind-représentables sur $C$), il faut et il suffit que $\varphi_\alpha$ soit} cofinal\textit{, ou ce qui revient au même, $I$ étant filtrante, qu'il satisfasse les deux conditions habituelles :}
\begin{itemize}[leftmargin=2.5em]
\item[F\,1)] $\forall\, \xi \in \Ob C/F$, $\exists\, i \in \Ob I$ \textit{et une flèche} $\xi \longrightarrow \varphi_\alpha(i)$.
\item[F\,2)] $\forall\, \xi \in \Ob C/F$, \textit{et} $i \in \Ob I$, \textit{et toute double-flèche}

\msp{page 80}
\[
\xi \rightrightarrows \varphi_\alpha(i) \;,
\]
\textit{il existe une flèche $i \xrightarrow{\;u\;} j$ dans $I$ telle que $\varphi_\alpha(u)$ égalise la double-flèche.}
\end{itemize}
\textsc{Démonstration}. Si $\varphi_\alpha$ est cofinal, on en conclut par définition que $\varinjlim \varphi_\alpha \longrightarrow \varinjlim \psi\varphi_\alpha$ est un isomorphisme, or cette flèche n'est autre que $\alpha$. Inversement, supposons que cette flèche soit un isomorphisme, donc $C/F \xleftarrow{\;\sim\;} C/\varinjlim x_i$, et il reste à voir que le foncteur canonique
\[
I \;\longrightarrow\; C/\varinjlim{}_{C^{\wedge}} x_i
\]
est cofinal, i.e. satisfait F\,1) et F\,2). La vérification est $\pm$ tautologique, en utilisant le fait que pour tout $y$ dans $C$, on a
\[
\Hom(y,\, \varinjlim{}_{C^{\wedge}} x_i) \;\xleftarrow{\;\sim\;}\; \varinjlim{}_{\mathrm{Ens}}\Hom(y,x_i),
\]
et du “calcul” des limites inductives filtrantes dans Ens.

% margin: vertical bar in the left margin beside Corollaire 4.4.3
\textbf{Corollaire 4.4.3.} \textit{Soit $\pi$ un ordinal infini. Pour que $F$ soit dans} $\Ind'_\pi(C)$ \textit{(image essentielle de} $\Ind_\pi(C)$ \textit{dans $C^{\wedge}$), il faut et il suffit qu'il existe une catégorie ordonnée grande devant $\pi$ $I$, et un foncteur cofinal $I \longrightarrow C/F$.}

Il nous faut donc donner des conditions, sur une catégorie $J$ $(= C/F)$, pour qu'il existe une catégorie ordonnée $I$ grande devant $\pi$, et un foncteur cofinal $I \longrightarrow J$.

\msp{page 81}
Une telle catégorie s'appellera \textit{grande devant} $\pi$. (Il est immédiat que dans le cas où $J$ est elle-même ordonnée, cette terminologie est cohérente avec celle déjà adoptée pour les ensembles ordonnés, i.e. que s'il existe une application croissante cofinale $\varphi : I \longrightarrow J$, avec $I$ grand devant $\pi$, alors $J$ est lui-même grand devant $\pi$.) Une telle catégorie est
\printedpage{17}
% head: A. Grothendieck : Dérivateurs, Ch. XVIII, Édition des Universités de Montpellier II et Paris VII
automatiquement filtrante (SGA 4 I 8.1.3 b)). Notons aussi que si on a un foncteur cofinal $J \longrightarrow J'$, si $J$ est grande devant $\pi$, $J'$ l'est aussi (par transitivité de la cofinalité).

\textbf{Proposition 4.4.4.} \textit{Soient $J$ une $\mathfrak{U}$-catégorie, $\pi$ un cardinal infini. Conditions équivalentes :}
\begin{itemize}[leftmargin=4em]
\item[\textit{a)}] \textit{$J$ est grande devant $\pi$, i.e. il existe un ensemble ordonné grand devant $\pi$ et un foncteur cofinal $\varphi : I \longrightarrow J$.}
\item[\textit{b)}] \textit{$J$ satisfait les deux conditions suivantes :}
\begin{itemize}[leftmargin=4em]
\item[$\mathrm{PF}_{\pi}\,1$] $\mathrm{Ord}\,J$ \textit{est grande devant $\pi$, i.e. pour toute famille $(y_\alpha)_{\alpha\in A}$ d'objets de $J$, avec} $\card A \leqslant \pi$\textit{, il existe un majorant $y$ dans $J$ (i.e. un $y$ et des flèches $y_\alpha \xrightarrow{\;v_\alpha\;} y$).}
\item[$\mathrm{PF}_{\pi}\,2$] \textit{Pour $y_0$, $y_1$ dans $J$ et tout ensemble $(u_\alpha)_{\alpha\in A}$ de flèches $y_0 \xrightarrow{\;u_\alpha\;} y_1$ avec} $\card A \leqslant \pi$\textit{, il existe une flèche $y_1 \xrightarrow{\;u'\;} y_2$ qui les égalise, i.e. telle que $u'u_\alpha = u'u_\beta \;\forall\, \alpha,\beta \in A$. (Il suffit de l'exiger quand} $\card A \geqslant 2$.\textit{)}
\end{itemize}
\item[\textit{c)}] \textit{Pour toute partie $A$ de} $\Fl J$\textit{, avec} $\card A \leqslant \pi$\textit{, existe une sous-catégorie $J'$ de $J$, avec} $\card\Fl J' \leqslant \pi$\textit{, telle que $A \subset \Fl J'$ et que $J'$ admette un objet final.}
\item[\textit{d)}] \textit{Soit $I$ l'ensemble des sous-catégories $J'$ de $J$ telles que} $\card J' \leqslant \pi$ \textit{et ayant un objet final}

\msp{page 82}
\textit{et ordonnons $I$ par inclusion. Alors $I$ est grand devant $\pi$, de plus le foncteur canonique $I \longrightarrow J$ $(J' \longmapsto e_{J'} =$ objet final de $J')$ est cofinal} (\footnote{Introduire $I(J)$ avant l'énoncé de la proposition. \textbf{N.B.} La condition de cofinalité $I(J) \longrightarrow J$ n'est pas automatique, comme on voit en prenant $J = B_G$, $G = \{\pm 1\}$.},\footnote{\textbf{N.B.} L'hypothèse de cofinalité est superflue si $J$ est rigide, i.e. tout endomorphisme d'un objet de $J$ est identique.},\footnote{\ N. Éd.\ \ Dans cette condition (d), le “\textit{de plus le foncteur \dots\ est cofinal}” est ajouté par Grothendieck \textit{a posteriori}, ce qui explique sa vaine tentative de l'établir dans l'implication d $\Longrightarrow$ a ci-dessous, implication qui devient évidente après cet ajout. Par ailleurs, les conditions (c) et (d) ne sont \textit{pas} impliquées par les conditions équivalentes (a) et (b). Par exemple, si $J$ est la catégorie ayant un seul objet et un seul morphisme non identique $p$ tel que $p^2 = p$, alors $J$ satisfait aux conditions (a) et (b), mais \textit{pas} aux conditions (c) et (d).}).
\end{itemize}
(\footnote{Autres conditions :
\begin{itemize}[leftmargin=3em]
\item[e)] Pour tout ensemble $A$ avec $\card A \leqslant \pi$, $\mathrm{diag}_J : J \longrightarrow J^A$ est cofinal.
\item[e$'$)] (Si $A_0$ est un ensemble donné, de cardinal $= \pi$.) $J \longrightarrow J^{A_0}$ est cofinal.
\end{itemize}
Voir aussi 4.9.6.9.1 et 4.9.6.9.2, pages 237, 238.})

\textsc{Démonstration}.
\begin{itemize}[leftmargin=7em,labelwidth=6em,align=left]
\item[a $\;\Longrightarrow\;$ b] vérification immédiate, en utilisant le critère F 1, F 2 de la cofinalité pour $I$ filtrante.
\item[b $\;\Longrightarrow\;$ c] par AQT \texttt{[âne qui trotte]}.
\item[c $\;\Longrightarrow\;$ d] par AQT \texttt{[âne qui trotte]}.
\end{itemize}
% foot: Version 17 décembre 2025

\printedpage{18}
% head: A. Grothendieck : Dérivateurs, Ch. XVIII, Édition des Universités de Montpellier II et Paris VII
\begin{itemize}[leftmargin=7em,labelwidth=6em,align=left]
\item[d $\;\Longrightarrow\;$ a] Considérons l'ensemble ordonné grand devant $\pi$ construit dans l'énoncé de d), et le \textit{foncteur}
\[
I \;\longrightarrow\; J
\]
donné par $J' \longmapsto e_{J'}$ $=$ objet final de $J'$. (Vérification que c'est un foncteur immédiate.) Ce foncteur est surjectif, donc satisfait la condition F 1 de cofinalité. La condition F 2 résulte aussi de d) en regardant les trois objets $J'$, $j \xrightarrow{\;u_1\;} e_{J'}$, $j \xrightarrow{\;u_2\;} e_{J'}$ de $I$,
\[
\begin{tikzcd}[column sep=small, row sep=small]
J', & e_{J'} \\
j \arrow[ur, shift left, "u_1"] \arrow[ur, shift right, "u_2"'] &
\end{tikzcd}
\]
et en prenant un $J''$ qui les majore tous trois, alors $e_{J'} \longrightarrow e_{J''}$ égalise $u_1$, $u_2$ (\footnote{Un peu bref, le cas où $e_{J'} = j$ échappe à cet argument, et demande qu'on \texttt{[phrase incomplète]}}).
\end{itemize}

% margin: vertical bar in the left margin beside Corollaire 4.4.5
\textbf{Corollaire 4.4.5.} \textit{Soit $F \in \Ob C^{\wedge}$. Pour que $F$ appartienne à} $\Ind'_\pi(C)$\textit{, il faut et il suffit que $C/F$ satisfasse à l'ensemble des quatre conditions équivalentes de 4.4.4.}

\textbf{Corollaire 4.4.6.} \textit{Soit $J$ une catégorie stable par} $\varinjlim$ \textit{de cardinal $\leqslant \pi$. Alors $J$ est grande devant $\pi$.}

En effet, les conditions $\mathrm{PF}_{\pi}\,1$, $\mathrm{PF}_{\pi}\,2$ sont visiblement vérifiées.

\msp{page 83}
\textbf{4.4.6.} Je me donne un foncteur % [sic] 4.4.6. -> 4.4.7. (?), numéro déjà employé par le Corollaire 4.4.6 qui précède
\[
J \xrightarrow{\;\varphi\;} \Ind'_\pi C \;, \tag*{(4.4.6.1)}
\]
$J$ une petite catégorie, et on cherche une $\varinjlim$ dans $\Ind'_\pi C$, en postulant au besoin l'existence des $\varinjlim$ de cardinal $\leqslant \pi$ dans $C$, et l'hypothèse $\card\Fl J \leqslant \pi$. Le procédé le plus naturel pour ramener la question de l'existence des $\varinjlim$ dans $\Ind'_\pi C$, à la question similaire dans $C$, est explicité dans SGA 4 I 8.8.3. On considère l'inclusion pleinement fidèle
\[
C \;\hookxrightarrow{\;\alpha\;}\; \Ind'_\pi C \;, \tag*{(4.4.6.2)}
\]
laquelle commute aux $\varinjlim$ de cardinal $\leqslant \pi$ (celles qui sont représentables dans $C$, cf. p. 78 et référence à SGA 4 I 9.8), d'où
\[
\uHom(J,C) \;\hookxrightarrow{\;\alpha\;}\; \uHom(J,\Ind'_\pi C) \tag*{(4.4.6.3)}
\]
également pleinement fidèle, et commutant aux $\varinjlim$ de cardinal $\leqslant \pi$ (représentables dans le premier membre). Cette flèche s'insère dans le carré commutatif relatif à la donnée d'un $j \in \Ob J$, et des foncteurs d'inclusion
\[
e \xrightarrow{\;\varepsilon_j\;} J \;, \qquad e_j(e) \;=\; j \;; \tag*{(4.4.6.4)} % [sic] e_j(e) -> \varepsilon_j(e), la flèche est notée \varepsilon_j
\]
% foot: Version 17 décembre 2025

% ---- from der-ch18_p021-025.tex
\printedpage{19}
% head: A. Grothendieck : Dérivateurs, Ch. XVIII, Édition des Universités de Montpellier II et Paris VII
% foot: Version 17 décembre 2025

\msp{page 84}
\[
\begin{tikzcd}
\uHom(J,C) \arrow[r, "\alpha^J"] \arrow[d, "\varepsilon_j^*"] & \uHom(J,\Ind'_\pi C) \arrow[d, "{\varepsilon_j^{*\prime}}"] & \ni \varphi \arrow[d, mapsto] \\
C \arrow[r, "\alpha"] & \Ind'_\pi(C) & \ni \varphi(j)\ ,
\end{tikzcd}
\qquad (\text{\footnotemark[21]})
\tag*{(4.4.6.2)}
\]
\footnotetext[21]{Les foncteurs $\varepsilon_j^*$, $\varepsilon_j^{*\prime}$ sont les foncteurs évaluation en $j$.}
lequel induit un foncteur
\[
\underbrace{\uHom(J/C)/\varphi}_{\overset{\text{déf}}{=}\, I} \xrightarrow{\varepsilon_j^\varphi} C/\varphi(j)\ . \tag*{(4.4.6.3)} % [sic] \uHom(J/C) -> \uHom(J,C), everywhere else written Hom(J,C)
\]
Si $C$ est stable par $\varinjlim$ relatives à un ensemble $K$ du type de conditions $\leqslant \pi$ (p.ex. $\varinjlim$ finies), il en est de même de $\uHom(J,C)$, et aussi de $\uHom(J,C)/\varphi$, et de $C/\varphi(j)$, en vertu du fait que $\alpha^J$ et $\alpha$ commutent aux $\varinjlim$ du type envisagé, de plus (4.4.6.3) commute aux limites en question. Supposons donc $C$ stable tout au moins par $\varinjlim$ finies, donc les catégories source et but dans (4.4.6.3) sont filtrantes. (Le fait que le foncteur commute aux $\varinjlim$ finies ne semble pas utile dans le contexte présent.) Cela nous servira pour avoir le critère commode de cofinalité SGA 4 I 8.1.3 b) par F\,1, F\,2.

\msp{page 85}
Rappelons que la cofinalité des foncteurs (4.4.6.3), quand elle a lieu ( ?), implique le fait que $\varphi$ est une limite inductive du foncteur composé
\[
\begin{tikzcd}[row sep=small]
I \arrow[r, "\psi_J"] \arrow[d, equal] & \uHom(J,C) \arrow[r, "\alpha^J"] \arrow[dd, dotted, "\varepsilon_j^*"] & \uHom(J,\Ind'_\pi(C)) \arrow[dd, dotted, "{\varepsilon_j^{*\prime}}"] \\
\uHom(J,C)/\varphi \arrow[d, "\varepsilon_j^\varphi"'] & & \\
C/\varphi(j) \arrow[r, dotted, "\psi_j"] & C \arrow[r, dotted, "\alpha"] & \Ind'_\pi(C)\ ,
\end{tikzcd}
\tag*{(4.4.6.4)}
\]
suivant l'ensemble ordonné filtrant $I$, plus précisément que $\varphi$ est la $\varinjlim$ “argument par argument en $j$”, ce qui signifie que pour tout $j \in \Ob J$, $\varphi(j) = \varepsilon_j^*(\varphi)$ est limite inductive de $\varepsilon_j^*(\alpha^J\psi) = (\alpha\psi_j)\varepsilon_j^\varphi$. (Or on sait en effet que $\varphi(j)$ est la limite de $\alpha\psi$, puisque $\varphi(j) \in \Ob C^\wedge$ est ind-représentable, donc aussi celle de $(\alpha\psi_j)\varepsilon_j^\varphi$ si $\varepsilon_j^\varphi$ est cofinal.)

Ainsi, admettant la cofinalité de (4.4.6.3) et interprétant $\psi_J$ dans (4.4.6.4)

\msp{page 86}
\[
\psi_J \ :\ I \longrightarrow \uHom(J,C)
\]
comme un foncteur
\[
I \times J \xrightarrow{\Psi} C\ , % [sic] \Psi -> \psi, the functor I x J -> C is called \psi in the next formula and in (4.4.7.1)
\]
on trouve la “représentation indicielle” de $\varphi$
\[
\varphi(j) \ =\ \underbrace{\varinjlim_{i\in I} \psi(i,j)}_{\substack{\text{limite inductive}\\ \text{prise dans } \Ind'_\pi(C)\\ \text{d'objets de } C \hookrightarrow \Ind'_\pi(C)}} \qquad ,
\]

\printedpage{20}
% head: A. Grothendieck : Dérivateurs, Ch. XVIII, Édition des Universités de Montpellier II et Paris VII
% foot: Version 17 décembre 2025
d'où
\[
\begin{aligned}
\underbrace{\varinjlim_{j\in J}{}_{\,\Ind'_\pi C}\varphi(j)}_{\text{sous réserve d'existence}} \quad &= \quad \varinjlim_{j\in J}\, \varinjlim_{i\in I} \psi(i,j)\\
&= \quad \varinjlim_{i\in I}\, \varinjlim_{j\in J} \psi(i,j)\ ,
\end{aligned}
\]
toutes les limites inductives étant prises dans $\Ind'_\pi(C)$. Mais si les $\varinjlim$ de type $J$ sont représentables dans $C$, et $\card J \leqslant \pi$, on sait que la $\varinjlim\limits_{j\in J} \psi(i,j)$ \textit{dans} $C$ est aussi leur $\varinjlim$ dans $\Ind'_\pi C$, et on trouve donc que $\varinjlim\limits_{j\in J} \varphi(j)$ existe dans $\Ind'_\pi(C)$ si et seulement si $\varinjlim\limits_{i\in I} \Psi(i)$ existe dans $\Ind'_\pi(C)$, où pour $i \in I$, on pose
\[
\Psi(i) \ =\ \varinjlim_{j\in J} \psi(i,j) \qquad \textit{dans } C.
\]

\msp{page 87}
Résumons ce que nous avons obtenu :

\textbf{Lemme 4.4.7.} \textit{Soit $C$ une petite catégorie, $\pi$ un cardinal infini, $J$ une catégorie telle que $\card \Fl J \leqslant \pi$, $\varphi : J \longrightarrow \Ind'_\pi(C)$ un foncteur. Supposons les $\varinjlim$ de type $J$ existant dans $C$, et que les foncteurs $\varepsilon_j^\varphi$ ($j \in \Ob J$) de (4.4.6.3) soient \textup{cofinaux} (ce qui s'explicite de façon simple surtout si source et but sont filtrants, chose acquise si $C$ est stable par $\varinjlim$ finies). Alors $\varphi$ a une “représentation indicielle canonique”}
\[
\psi \ :\ I \times J \longrightarrow C \tag*{(4.4.7.1)}
\]
\textit{avec}
\[
I \ =\ \uHom(J,C)/\varphi\ , \tag*{(4.4.7.2)}
\]
\textit{i.e. on a l'isomorphisme fonctoriel en $j \in \Ob J$}
\[
\varphi(j) \ \simeq\ \varinjlim_{I} \psi(i,j) \qquad (\text{\footnotemark[22]})\ . \tag*{(4.4.7.3)}
\]
\footnotetext[22]{limite inductive dans $\Ind'_\pi$ d'objets de $C$.}
\textit{Posons pour tout $i \in \Ob I$}
\[
\Psi(i) \ =\ \varinjlim_{J} \psi(i,j) \qquad (\text{\footnotemark[23]})\ . \tag*{(4.4.7.4)}
\]
\footnotetext[23]{limite inductive prise dans $C$, mais valable aussi dans $\Ind'_\pi C$.}
\textit{Alors $\varinjlim\limits_{j\in J} \varphi(j)$ (dans $\Ind'_\pi C$) existe si et seulement si}
\msp{page 88}
\textit{la $\varinjlim\limits_{i\in I} \Psi(i)$ existe (dans $\Ind'_\pi C$ également), et les deux sont égales.}

\textbf{Corollaire 4.4.8.} \textit{Supposons que $C$ soit stable par $\varinjlim$ de cardinal $\leqslant \pi$ (pas seulement par celles de type $J$). Alors $\uHom(J,C)$ et $I = \uHom(J,C)/\varphi$ ont la même propriété, $I$ est grande devant $\pi$, et $\varinjlim\limits_{j\in J} \varphi(j) \ =\ \varinjlim\limits_{i\in I} \Psi(i)$ existe dans $\Ind'_\pi(C)$, sous réserve de cofinalité des foncteurs (4.4.6.3).}

\printedpage{21}
% head: A. Grothendieck : Dérivateurs, Ch. XVIII, Édition des Universités de Montpellier II et Paris VII
% foot: Version 17 décembre 2025
Il suffit en effet de vérifier ce

\textbf{Lemme 4.4.8.1.} \textit{Soit $I$ une catégorie grande devant $\pi$, et $\Psi : I \longrightarrow C$ un foncteur. Alors $\varinjlim\limits_{i} \Psi(i)$ existe dans $\Ind'_\pi(C)$, et c'est la limite inductive dans $C^\wedge$, ou encore “c'est” l'ind-objet grand devant $\pi$ représenté par $\Psi$.} (En d'autres termes (\footnote[24]{C'est abusif, cette stabilité est un énoncé nettement plus fort! Je le prouve page 99.}), $\Ind'_\pi(C) \subset C^\wedge$ est une sous-catégorie strictement pleine \textit{stable} par $\varinjlim$ grandes devant $\pi$.)

\msp{page 89}
Ce lemme est une tautologie : si $F$ désigne la limite inductive dans $C^\wedge$, il est clair par définition (en prenant un foncteur $I_0 \longrightarrow I$ cofinal, avec $I_0$ \textit{ordonnée} grande devant $\pi$, de sorte que la $\varinjlim$ peut être prise suivant $I_0$) que $F$ est dans $\Ind'_\pi(C)$, et alors c'est a fortiori une $\varinjlim$ dans $\Ind'_\pi C$.

Pour vérifier l'existence des $\varinjlim$ de cardinal $\leqslant \pi$ dans $\Ind'_\pi C$, dans le cas où $C$ lui-même est stable par lesdites limites, il suffit donc, grâce au corollaire 4.4.8, de prouver la cofinalité des foncteurs (4.4.6.3). Cela va résulter du

\textbf{Lemme 4.4.9.} \textit{Soient $C$ une petite catégorie, $\pi$ un cardinal infini, $J$ une petite catégorie telle que pour tous $j, j' \in \Ob J$, on ait $\card \Hom(j,j') \leqslant \pi$. Sous ces conditions, pour tout $j \in \Ob J$, le foncteur (4.4.6.3) \textup{admet un adjoint à gauche,} donc a fortiori il est cofinal (car les catégories colocalisées au dessus de $C/\varphi(j)$ ont un objet initial, a fortiori elles sont $0$-connexes).}

\msp{page 90}
\textsc{Démonstration de} 4.4.8. Revenons au diagramme (4.4.6.2), dont on déduit (4.4.6.3), nous allons le réécrire de façon simplifiée :
\[
\begin{tikzcd}
B \arrow[r, "\beta"] \arrow[d, shift left, "\varepsilon"] & B' \arrow[d, shift left, "\varepsilon'"] & \ni y'_0 \arrow[d, mapsto] \\
C \arrow[r, "\gamma"] \arrow[u, dotted, shift left, "\lambda"] & C' \arrow[u, dotted, shift left, "\lambda'"] & \ni z'_0\ .
\end{tikzcd}
\tag*{(4.4.9.1)}
\]
Je dis que $\varepsilon$, $\varepsilon'$ ont \textit{des adjoints à gauche} $\lambda$ resp. $\lambda'$, et que l'homomorphisme canonique de “changement de base”
\[
\lambda'\gamma \longrightarrow \beta\lambda \tag*{(4.4.9.2)}
\]
est un \textit{isomorphisme.} En effet, l'existence des adjoints de $\varepsilon$, $\varepsilon'$, à savoir des foncteurs $\varepsilon_{j\,!}$ et $\varepsilon'_{j\,!}$ relatifs à l'inclusion
\[
\varepsilon_j \ :\ e \longrightarrow J
\]
et aux catégories de foncteurs $\uHom(J,C)$ resp. $\uHom(J,\Ind'_\pi C)$, résultant de l'existence des sommes de type $\Hom(j,j')$ ($j' \in \Ob J$) dans $C$ (déjà acquise), \textit{et dans} $\Ind'_\pi C$\textit{ ;} et le fait que le foncteur $C \longrightarrow \Ind'_\pi C$ commute
\msp{page 91}
auxdites sommes (chose également déjà acquise, puisque $\card\, (\Hom(j,j')) \leqslant \pi$) implique alors que (4.4.4.2) est un isomorphisme. Moyennant la vérification en suspens (l'existence % [sic] (4.4.4.2) -> (4.4.9.2), the isomorphism in question is the change-of-base arrow (4.4.9.2)
\printedpage{22}
% head: A. Grothendieck : Dérivateurs, Ch. XVIII, Édition des Universités de Montpellier II et Paris VII
% foot: Version 17 décembre 2025
des sommes de cardinal $\leqslant \pi$ dans $\Ind'_\pi(C)$), le lemme 4.4.9 résultera du résultat plus général :

\textbf{Lemme 4.4.10.} \textit{Considérons un carré commutatif (4.4.9.1) (traits pleins) de catégories, où $\varepsilon$ admet un adjoint à gauche $\lambda$, $\varepsilon'$ un adjoint à gauche $\lambda'$, et la flèche de changement de base (4.4.9.2) est un isomorphisme. Alors pour tout $y'_0$ dans $B'$, posant $\varepsilon'(y'_0) = z'_0$, le foncteur induit par $(\varepsilon, \varepsilon')$,}
\[
\overline{B} \ =\ B/y'_0 \xrightarrow{\bar{\varepsilon}} \overline{C} = C/z'_0\ ,
\]
\textit{admet un adjoint à gauche $\overline{\lambda}$, donné par}
\[
\begin{tikzcd}[column sep=small]
\overline{\lambda}(\gamma(z) \xrightarrow{v} z'_0) \ \ = & \Bigl(\lambda'\gamma(z) \simeq \beta(\lambda(z)) \arrow[r, "{\lambda'(v)}"] \arrow[ddr, "{\substack{\overline{\lambda}(v) \,=\, \text{composé}\\ \text{adj} \circ \lambda'(v)}}"'] & \lambda'(z'_0) \arrow[d, equal] \\
 & & \lambda'\varepsilon'(y'_0) \arrow[d, "\text{adj.}"] \\
 & & y'_0\Bigr)\ \ .
\end{tikzcd}
\]

\msp{page 92}
La démonstration se fait par un AQT \texttt{[âne qui trotte]} fastidieux (j'avoue ne pas avoir poussé jusqu'au dernier détail \ldots).

Pour terminer la démonstration de 4.4.9, il reste à prouver un dernier ( ?)

\textbf{Lemme 4.4.11.} \textit{Soit $C$ une petite catégorie, stable par sommes de cardinal $\leqslant \pi$. Alors $\Ind'_\pi C$ est stable par sommes de cardinal $\leqslant \pi$ (et, pour mémoire, le foncteur $\alpha : C \longrightarrow \Ind'_\pi C$ y commute).}

Appliquant le corollaire 4.4.8 au lemme 4.4.7, on est ramené à prouver la cofinalité de (4.4.6.3) quand $J$ est discrète, plus le fait que $\uHom(J,C)/\varphi$ est grande devant $\pi$, dans ce cas. Mais la donnée de $\varphi$ revient à celle d'une famille $(\varphi_j)_{j\in J}$ d'objets de $\Ind'_\pi(C)$, et le foncteur (4.4.6.3) s'explicite comme la projection
\[
\prod_{j\in J} (C/\varphi_j) \longrightarrow C/\varphi_j
\]
du produit total sur un des facteurs. Donc on est ramené à ceci :

\msp{page 93}
\textbf{Lemme 4.4.12.} \textit{Soit $(C_j)_{j\in J}$ une famille de petites catégories grandes devant $\pi$, indexée par un petit ensemble $J$. Alors $\prod C_j$ est grande devant $\pi$, et les projections (pour $j_0 \in J$) $\prod C_j \longrightarrow C_{j_0}$ sont cofinales.}

Prouvons d'abord la cofinalité, qui est le plus important pour nous. C'est évident si $J = \{j_0\}$, sinon on a $J = J' \cup \{j_0\}$, et $\prod_{j\in J} C_j \simeq C' \times C_{j_0}$, où $C' = \prod_{j\in J'} C_j$, donc il faut voir que $C' \times C_{j_0} \longrightarrow C_{j_0}$ est cofinal, i.e. que toute catégorie $\xi\backslash(C' \times C_{j_0})$ (pour $\xi \in \Ob C_{j_0}$) est 0-connexe, or elle est isomorphe à $C' \times (\xi\backslash C_{j_0})$, et $\xi\backslash C_{j_0}$ étant 0-connexe, il reste à voir que $C'$ est 0-connexe. On fera attention à cet égard qu'un produit \textit{infini} de catégories 0-connexes n'est pas forcément 0-connexe, donc la cofinalité n'est pas totalement tautologique! Mais si on prouve que la catégorie $\prod_{j\in J} C_j$ est
\printedpage{23}
% head: A. Grothendieck : Dérivateurs, Ch. XVIII, Édition des Universités de Montpellier II et Paris VII
% foot: Version 17 décembre 2025
\msp{page 94}
grande devant $\pi$, donc filtrante, cela implique qu'elle est 0-connexe, et on peut appliquer ce même résultat à $C' = \prod_{j\in J'} C_j$, et on gagne.

Donc on est ramené à prouver que $\prod C_j$ est grande devant $\pi$. Ceci se voit soit sur le critère $\mathrm{PF}_\pi\,1$, $\mathrm{PF}_\pi\,2$ de 4.4.4 b, soit en se ramenant au cas où les $C_j$ sont ordonnées, mais cela a l'air plus long. Notons que l'on prouve de même qu'un (petit) produit de catégories filtrantes est encore une catégorie filtrante.

\textbf{Remarques.}
\begin{enumerate}[label={}, leftmargin=*]
\item[1.] Cet argument ne montre nullement (même si $C$ est stable par $\varinjlim$ de cardinal $\leqslant \pi$) que pour tout $\varphi \in \uHom(J,\Ind(C))$, le foncteur similaire à (4.4.6.3)
\[
\uHom(J,C)/\varphi \longrightarrow C/\varphi(j)
\]
soit cofinal, et je doute qu'il le soit toujours. La difficulté provient du fait que $C \longrightarrow \Ind\, C$ ne commute pas aux sommes de cardinal $\leqslant \pi$, seulement aux sommes finies,
\msp{page 95}
donc le lemme 4.4.10 ne peut être invoqué.
\item[2.] Nous venons de prouver, via 4.4.8 et 4.4.9 (dont la démonstration est achevée avec 4.4.12), que si $C$ est une petite catégorie stable par $\varinjlim$ de cardinal $\leqslant \pi$, alors $\Ind_\pi C$ est également stable par ces limites. De plus, les systèmes inductifs de ce type $J$ dans $\Ind_\pi C$ admettent une représentation indicielle
\[
\psi \ :\ I \times J \longrightarrow C\ ,
\]
avec $I$ grande devant $\pi$. Et on peut toujours supposer, quitte à prendre $I' \longrightarrow I$ cofinal avec $I'$ ensemble ordonné grand devant $\pi$, que $I$ est de plus un ensemble ordonné.
\end{enumerate}

C'est là le moyen donc, pour la démonstration du

\textbf{Théorème 4.4.13.} \textit{Soit $C$ une catégorie équivalente à une petite catégorie, et stable par $\varinjlim$ de cardinal $\leqslant \pi$, où $\pi$ est un cardinal infini donné. Alors :}
\msp{page 96}

\begin{enumerate}[label={}, leftmargin=*]
\item[\textit{a)}] \textit{(Pour mémoire.) Le foncteur canonique}
\[
C \longrightarrow \Ind_\pi(C) \xrightarrow{\ \sim\ } \Ind'_\pi(C)
\]
\textit{commute aux $\varinjlim$ de cardinal $\leqslant \pi$.}
\item[\textit{b)}] \textit{(Pour mémoire.) Les systèmes inductifs de cardinal $\leqslant \pi$ dans $\Ind_\pi(C)$ (i.e. les foncteurs $\varphi : J \longrightarrow \Ind_\pi(C)$, avec $\card \Fl(J) \leqslant \pi$) admettent une représentation indicielle}
\[
\psi \ :\ I \times J \longrightarrow C\ ,
\]
\textit{avec $I$ ensemble ordonné grand devant $\pi$, et admettent une limite inductive dans $\Ind_\pi C$ qui se calculent par}
\[
\underbrace{\varinjlim_{j\in J} \varphi(j)}_{\text{dans } \Ind_\pi(C)} \ \simeq\ \text{“}\,\varinjlim_{i\in I}\,\text{”}\Psi(i)
\]
\end{enumerate}

% ---- from der-ch18_p026-030.tex
\printedpage{24}
% head: A. Grothendieck : Dérivateurs, Ch. XVIII, Édition des Universités de Montpellier II et Paris VII
% foot: Version 17 décembre 2025

\textit{avec}
\[
\Psi(i) \;=\; \varinjlim_{j\in J}{}_{C}\,\psi(i,j)\;.
\]
\textit{c)} $\Ind_\pi C$ \textit{est stable par petites limites inductives} (\footnote{rajouter aussi, en d), le corollaire 4.4.16, page 101.}).

Il reste à établir ce dernier point. Comme on a déjà les $\varinjlim$ finies par b, il reste à prouver l'existence des petites $\varinjlim$ filtrantes, ou au choix des sommes infinies, ce qui sera plus joli en fait. Mais
\msp{page 97}
on pourra aussi bien traiter le cas général d'une $\varinjlim$ de type $J$ quelconque ($J$ petite). Pour ceci, soit
\[
I \;=\; I(J) \;=\; \begin{aligned}[t]
&\text{ensemble des sous-catégories } J' \text{ de } J \text{ telles que}\\
&\card\Fl(J')\leqslant\pi,
\end{aligned}
\]
ordonnons $I$ par inclusion. Il est évident que $I\neq\emptyset$ (car $\emptyset\in I$), et même que $I$ \textit{est grand devant} $\pi$, car si $(J'_\alpha)_{\alpha\in A}$ avec $\card(A)\leqslant\pi$ est une famille d'éléments de $I$, alors la sous-catégorie de $J$ engendrée par les $J'_\alpha$ est encore dans $I$ (vérification immédiate), et majore les $J'_\alpha$. En fait, la considération de $I(J)$ associée à $J$ va nous permettre de prouver ceci, qui généralise 4.4.13 c), compte tenu de b) :

\textbf{Corollaire 4.4.14} (\footnote{Comparer aussi 4.11.3, page 270.}). \textit{Soit $\mathcal{M}$ une catégorie, $\pi$ un cardinal infini tel que}
\begin{enumerate}
\item[a)] \textit{$\mathcal{M}$ est}
\msp{page 98}
\textit{stable par $\varinjlim$ de cardinal $\leqslant\pi$, et}
\item[b)] \textit{$\mathcal{M}$ est stable par $\varinjlim$ sur les ensembles ordonnés $I$ grands devant $\pi$.}
\end{enumerate}
\textit{Alors $\mathcal{M}$ est stable par petites $\varinjlim$ .}

(\footnote{\textbf{Corollaire 4.4.14.1.} $\mathcal{M}$ étant comme dans 4.4.14, et $f:\mathcal{M}\longrightarrow\mathcal{E}$ étant un foncteur qui commute aux $\varinjlim$ de cardinal $\leqslant\pi$, et aux petites $\varinjlim$ grandes devant $\pi$, alors $f$ commute aux petites $\varinjlim$ .}).

En effet, si $\varphi:J\longrightarrow\mathcal{M}$ est donné, on en construit
\[
\psi\;:\;I\;=\;I(J)\;\longrightarrow\;\mathcal{M}
\]
par
\[
\psi(J')\;=\;\varinjlim_{J'}\varphi(J')\;,
\]
qui existe par hypothèse sur $\mathcal{M}$. Il est alors immédiat que $\varinjlim\limits_{J}\varphi$ existe si et seulement si $\varinjlim\limits_{I}\psi$ existe, or cette existence est assurée du fait que $I$ est grand devant $\pi$, et l'hypothèse sur $\mathcal{M}$.

Pour achever de prouver le théorème 4.4.13, on est donc ramené au cas particulier des limites inductives $\varinjlim\limits_{J}\varphi(j)$, quand $\varphi:J\longrightarrow\Ind'_\pi C$ est donné, avec $J$ ensemble
\printedpage{25}
% head: A. Grothendieck : Dérivateurs, Ch. XVIII, Édition des Universités de Montpellier II et Paris VII
% foot: Version 17 décembre 2025
\msp{page 99}
ordonné grand devant $\pi$. Mais ceci va résulter du résultat plus précis :

\textbf{Corollaire 4.4.15} (\footnote{comparer à 4.7.3.1, qu'il y aurait lieu d'expliciter ici.}). \textit{Soit $C$ une petite catégorie. Alors la sous-catégorie strictement pleine} $\Ind'_\pi C$ \textit{de $C^\wedge$ est stable dans $C^\wedge$ par $\varinjlim$ grandes devant $\pi$, en d'autres termes, ces $\varinjlim$} existent dans $\Ind'_\pi(C)$, \textit{et le foncteur d'inclusion y commute.} % doubt resolved: existent is upright in the PDF; see corrections.csv, PDF leaf 27

\textsc{Démonstration}. Soit $I$ ensemble ordonné, grand devant $\pi$, et $i\longmapsto F_i$ un foncteur $I\longrightarrow\Ind'_\pi(C)$, $F\;=\;\varinjlim\limits_{I}F_i$ dans $C^\wedge$, prouvons que $F\in\Ob\Ind'_\pi(C)$, i.e. (4.4.3) que $C/F$ est grande devant $\pi$. J'utilise le critère b) de 4.4.4, il faut vérifier $\PF_\pi\,1$, $\PF_\pi\,2$. Soient
\[
x_\alpha\xrightarrow{\;u_\alpha\;}F\;=\;\varinjlim F_i
\]
des objets de $C/F$ indexés par $A$ avec $\card A\leqslant\pi$. Chaque $u_\alpha$ provient de $x_\alpha\longrightarrow F_{i_\alpha}$, et en prenant un majorant commun des $i_\alpha$ ($I$ grand devant $\pi$), on trouve un $i\in I$ et des $v_\alpha:x_\alpha\longrightarrow F_i$. Comme $F_i$ est dans $\Ind'_\pi(C)$, i.e. $C/F_i$ grande devant $\pi$, considérant les $v_\alpha$ comme des objets de $C/F_i$, ils sont majorés par un objet $x\longrightarrow F_i$ de $C/F_i$. On a donc
\[
\begin{tikzcd}
x_\alpha \arrow[r] \arrow[d] & F_{i_\alpha} \arrow[d] \arrow[dr, dotted] & \\
x \arrow[r] & F_i \arrow[r] & F\;,
\end{tikzcd}
\]
qui montre que les $x_\alpha\longrightarrow F$ sont majorés par $x\longrightarrow F$,
\msp{page 100}
ce qui prouve $\PF_\pi\,1$. Prouvons $\PF_\pi\,2$, i.e. soient $x\xrightarrow{\;u\;}F$, $y\xrightarrow{\;v\;}F$ deux objets de $C/F$, et $(f_\alpha)_{\alpha\in A}$ une famille de cardinal $\leqslant\pi$ de flèches de $u$ dans $v$. On sait que $u$, $v$ proviennent de $u_i:x\longrightarrow F_i$, $v_i:y\longrightarrow F_i$ pour $i$ assez grand.
\[
\begin{tikzcd}[column sep=small]
x \arrow[rr, "f_\alpha"] \arrow[dr, "u"'] & & y \arrow[dl, "v"] \\
& F &
\end{tikzcd}
\]
\[
\begin{tikzcd}[column sep=small]
x \arrow[rr, "f_\alpha"] \arrow[dr, "u_i"'] & & y \arrow[dl, "v_i"] \\
& F_i &
\end{tikzcd}
\]

Pour tout $\alpha$, on peut choisir $i_\alpha\geqslant i$ tel que $v_if_\alpha$ et $u_i$ deviennent égaux en composant avec $F_i\longrightarrow F_{i_\alpha}$, et comme $\card A\leqslant\pi$ et $I$ grand devant $\pi$, il existe un majorant $i'$ des $i_\alpha$. Au total, on peut supposer (remplaçant $i$ par $i'$) que $v_if_\alpha=u_i$ pour tout $\alpha$. Alors $u_i$, $v_i$ étant considérés comme objets de $C/F_i$, les $f_\alpha$ sont des flèches de $C/F_i$ de $u_i$ dans $v_i$. Comme $C/F_i$ est grand devant $\pi$ (par l'hypothèse $F_i\in\Ob\Ind'_\pi(C)$), on trouve qu'il existe un égalisateur $(y,v_i)\xrightarrow{\;g\;}(z,w_i)$.
\printedpage{26}
% head: A. Grothendieck : Dérivateurs, Ch. XVIII, Édition des Universités de Montpellier II et Paris VII
% foot: Version 17 décembre 2025
\[
\begin{tikzcd}
x \arrow[r, shift left=2, "(f_\alpha)"] \arrow[r] \arrow[r, shift right=2] \arrow[dr, "u_i"'] & y \arrow[r, "g"] \arrow[d, "v_i"'] & z \arrow[dl, "w_i"] \\
& F_i &
\end{tikzcd}
\]

Alors l'image de $g$ dans $C/F$ (via $F_i\longrightarrow F$) égalise les flèches $f_\alpha$, en tant que flèches de $C/F$, ce qui prouve $\PF_\pi\,2$ et achève la démonstration de 4.4.14, et par là aussi celle du théorème 4.4.13, et du lemme 4.3. Par là est aussi achevée la démonstration de la première assertion du lemme
\msp{page 101}
4.3.7 (p. 77), ce qui achève la démonstration du théorème 4.3.6 (p. 70-71). Il reste à établir la propriété universelle énoncée dans 4.3.7, que je vais rappeler :

\textbf{Corollaire 4.4.16} (\footnote{Pourrait être incorporé en partie d) au théorème 4.4.13, car c'est un complément important, et pas tellement évident finalement.}). \textit{Soient $C$ une catégorie équivalente à une petite catégorie, stable par $\varinjlim$ de cardinal $\leqslant\pi$, où $\pi$ est un cardinal infini donné, et considérons le foncteur d'inclusion canonique}
\[
C\xrightarrow{\;\alpha\;}\Ind_\pi(C)\;. \tag*{(4.4.16.1)}
\]
\textit{Ce foncteur est} 2\textit{-universel pour les foncteurs $f:C\longrightarrow\mathcal{E}$, où $\mathcal{E}$ est une $\mathcal{U}$-catégorie stable par petites $\varinjlim$ , et $f$ un foncteur commutant à ces limites inductives. En d'autres termes, on a}
\begin{enumerate}
\item[\textit{1$^\circ$)}] \textit{$\alpha$ possède les propriétés précédentes, i.e.} $\Ind_\pi(C)$ \textit{est stable par $\varinjlim$ , et le foncteur $\alpha$ commute aux $\varinjlim$ de cardinal $\leqslant\pi$.}
\item[\textit{2$^\circ$)}] \textit{Pour toute $\mathfrak{U}$-catégorie $\mathcal{E}$ stable par petites limites inductives, le foncteur}
\[
\begin{aligned}
g\;&\longmapsto\;g\circ\alpha\\
\uHom_{!}(\Ind_\pi C,\mathcal{E})\;&\longrightarrow\;\uHom_{!\,\pi}(C,\mathcal{E})
\end{aligned}
\tag*{(4.4.16.2)}
\]
\textit{est une équivalence de catégories, où l'indice $!$ (resp. $!\,\pi$) indique la sous-catégorie strictement pleine de la catégorie} $\uHom$
\msp{page 102}
\textit{envisagée, formée des foncteurs qui commutent aux petites $\varinjlim$ (resp. à celles de cardinal $\leqslant\pi$).}
\end{enumerate}

On peut donc dire aussi que $\Ind_\pi(C)$ joue le rôle d'une “enveloppe” stable par petites $\varinjlim$ de la catégorie $C$, quand les $\varinjlim$ de cardinal $\leqslant\pi$, lesquelles existent déjà dans $C$, sont respectées par les plongement de $C$ dans son enveloppe. C'est un résultat nullement tautologique qu'il suffit pour cela d'ajouter à $C$ des $\varinjlim$ \textit{grandes devant} $\pi$ d'objets de $C$, et que plus précisément, la catégorie enveloppe est réalisée par $\Ind_\pi C$ (ou, au choix et moins pléthoriquement, par $\Ind'_\pi C\subset C^\wedge$). % [sic] les plongement -> le plongement, number agreement

\textsc{Démonstration}. Le 1°) est contenu dans 4.4.13 a) et c). Prouvons donc 2°). Il est évident que $\alpha$ est un foncteur pleinement fidèle dont l'image essentielle engendre strictement
\printedpage{27}
% head: A. Grothendieck : Dérivateurs, Ch. XVIII, Édition des Universités de Montpellier II et Paris VII
% foot: Version 17 décembre 2025
$\mathcal{M}=\Ind_\pi C$, d'où résulte que le foncteur (4.4.16.2) est pleinement fidèle, étant induit par le foncteur pleinement fidèle
\[
\alpha^*_{\mathcal{E}}\;:\;\uHom(\mathcal{M},\mathcal{E})\;\longrightarrow\;\uHom(C,\mathcal{E})
\]
\msp{page 103}
sur les sous-catégorie pleines envisagées. Il reste donc à montrer que (4.4.16.2) est essentiellement surjectif. Soit donc $f:C\longrightarrow\mathcal{E}$ dans $\uHom_{!\,\pi}(C,\mathcal{E})$, considérons son extension canonique % [sic] les sous-catégorie pleines -> les sous-catégories pleines, number agreement
\[
\left\{
\begin{aligned}
&Lf\;:\;\Ind(C)\;\longrightarrow\;\mathcal{E}\\
&Lf((X_i)_{i\in I})\;=\;\varinjlim_{i\in I}f(X_i)
\end{aligned}
\right.
\tag*{(4.4.16.3)}
\]
(cf. SGA 4 I 8.7.3), et soit
\[
\psi(f)\;=\;Lf|\Ind_\pi(C)\;. \tag*{(4.4.16.4)}
\]
Rappelons que $Lf$ commute aux petites $\varinjlim$ filtrantes, mais comme le foncteur d'inclusion
\[
\Ind_\pi C\;\longrightarrow\;\Ind(C) \tag*{(4.4.16.5)}
\]
ne commute pas à ces $\varinjlim$ , il n'est pas clair a priori que $\psi(f)$ commute auxdites limites, ni a fortiori qu'il commute aux petites $\varinjlim$ sans plus. Mais on a vu que (4.4.16.5) commute aux $\varinjlim$ \textit{grandes devant} $\pi$, ce qui est contenu dans le corollaire 4.4.15 (vu que les limites inductives \textit{filtrantes} dans $\Ind(C)$ sont les limites \textit{dans} $C^\wedge$) : dire que $\Ind'_\pi(C)\longrightarrow\Ind'(C)$ commute à un certain type de $\varinjlim$ filtrantes, revient à dire qu'il en est ainsi du composé avec $\Ind'_\pi(C)\hookrightarrow C^\wedge$, ce qui est établi par 4.4.15. De ceci résulte que $\psi(f)$ \textit{commute aux} $\varinjlim$ \textit{grandes devant} $\pi$. Pour établir que $g=\psi(f)$ commute
\msp{page 104}
aux petites $\varinjlim$ , il reste donc à voir qu'il commute aux $\varinjlim$ de cardinal $\leqslant\pi$, en vertu de 4.4.14.1 (marge page 98). Pour l'établir, on utilise le théorème 4.4.13 b, en représentant un foncteur $\varphi:J\longrightarrow\Ind_\pi(C)$ (où $\card\Fl J\leqslant\pi$) sous forme indicielle par
\[
\Phi\;:\;I\times J\;\longrightarrow\;C\;, \tag*{(4.4.16.6)}
\]
$I$ étant un ensemble ordonné grand devant $\pi$. On a donc
\[
\varphi(j)\;=\;\varinjlim_{i\in I}{}_{\Ind_\pi C}\,\Phi(i,j)\;,
\]
et comme $g=\psi(f)$ commute aux $\varinjlim$ de type $I$ par ce qu'on a vu, on aura
\[
g(\varphi(j))\;=\;\varinjlim_{i\in I}{}_{(\mathcal{E})}\;g(\Phi(i,j))\;,
\]
d'où
\[
\varinjlim_{j\in J}{}_{(\mathcal{E})}\,g(\varphi(j))\;=\;\varinjlim_{j\in J}\;\varinjlim_{i\in I}\;\underbrace{g(\Phi(i,j))}_{=\,f(\Phi(i,j))}\qquad\text{(limites dans $\mathcal{E}$)}\;. \tag*{$(*)$}
\]
\printedpage{28}
% head: A. Grothendieck : Dérivateurs, Ch. XVIII, Édition des Universités de Montpellier II et Paris VII
% foot: Version 17 décembre 2025

D'autre part, on a
\[
\begin{aligned}
\varinjlim_{j\in J}{}_{\Ind_\pi}\,\varphi(j)\;&\simeq\;\varinjlim_{j\in J}\;\varinjlim_{i\in I}\;\Phi(i,j)\\
&\simeq\;\varinjlim_{i\in I}\;\underbrace{\varinjlim_{j\in J}\;\Phi(i,j)}_{\substack{\text{peut être calculé}\\ \text{dans $C$, soit}\\ \Psi(i)\,\in\,\Ob C}}\;,
\end{aligned}
\]
donc
\[
\varinjlim_{j\in J}{}_{\Ind_\pi}\,\varphi(j)\simeq\varinjlim_{i\in I}\;\Psi(i)\;, \tag*{$(**)$}
\]
\msp{page 105}
d'où en appliquant $g$, qui commute aux $\varinjlim$ de type $I$, et se rappelant maintenant que $g|C\simeq f$,
\[
g(\,\varinjlim_{J}{}_{\Ind_\pi C}\,\varphi(j))\;\simeq\;\varinjlim_{i\in I}\;f(\Psi(i))\;. \tag*{$(*{*}*)$}
\]
On a d'autre part
\[
\Psi(i)\;\overset{\text{déf}}{=}\;\varinjlim_{j\in J}\;\Phi(i,j)\;,\qquad\qquad \varinjlim\;\text{ dans } C\;,
\]
et comme $f$ commute aux $\varinjlim$ de type $J$ (car $\card\Fl J\leqslant\pi$), on en conclut
\[
f(\Psi(i))\;\simeq\;\varinjlim_{j\in J}\;f(\Phi(i,j))\;,
\]
donc par $(*{*}*)$
\[
g(\,\varinjlim_{j\in J}\;\varphi(j))\;\simeq\;\varinjlim_{i\in I}\;\varinjlim_{j\in J}\;f(\Phi(i,j))\;. \tag*{$(*{*}{*}*)$}
\]
En composant avec $(*)$ et utilisant à nouveau la commutativité de $\varinjlim\limits_{I}$ à $\varinjlim\limits_{J}$ , on trouve donc
\[
g(\,\varinjlim\;\varphi(j))\;\overset{\sim}{\longleftarrow}\;\varinjlim\;g(\varphi(j))\;,
\]
ce qui prouve que $g$ commute aux limites inductives de type $J$, et achève de prouver que $g$ est \texttt{[dans]} $\uHom_{!}(\Ind_\pi(C),\mathcal{E})$. Comme $g|C\simeq f$, on en conclut que le foncteur (4.4.16.2) est essentiellement surjectif, ce qui achève de prouver 4.4.16. De plus, on a construit un foncteur quasi-inverse explicite, $\psi$,
\msp{page 106}
s'insérant dans le diagramme commutatif
\[
\begin{tikzcd}[column sep=large, row sep=large]
\uHom_{!\,\pi}(C,\mathcal{E}) \arrow[r, "\psi", "(\sim)"'] \arrow[d, hook, "\text{incl. pl. fid.}"] & \uHom_{!}(\Ind_\pi(C),\mathcal{E}) \arrow[r] \arrow[d, hook, dotted, "\text{pl. fid.}"] & \uHom_{!\,\pi}(\Ind_\pi(C),\mathcal{E}) \\
\uHom(C,\mathcal{E}) \arrow[r, "L", "\substack{\sim\\ \text{(cf. SGA 4 I 8.7.3)}}"'] & \uHom_{!\,\mathrm{filt}}(\Ind(C),\mathcal{E}) \arrow[ur, "R\ \text{(restriction)}"'] &
\end{tikzcd}
\qquad.
\tag*{(4.4.16.6)} % [sic] (4.4.16.6) -> (4.4.16.7)?, the label (4.4.16.6) is already used for Phi on the preceding page
\]

% ---- from der-ch18_p031-035.tex
\setcounter{footnote}{29}

\printedpage{29}
% head: A. Grothendieck : Dérivateurs, Ch. XVIII, Édition des Universités de Montpellier II et Paris VII
% foot: Version 17 décembre 2025

\section*{4.5\quad Filtration cardinale canonique d'une catégorie $\mathfrak{U}$-paratopos}

\textbf{Définition 4.5.1.} On appelle \textit{catégorie $\mathfrak{U}$-paratopos} une catégorie $\mathcal{M}$ $\mathfrak{U}$-pseudo-topos et accessible, i.e. qui satisfait aux conditions équivalentes du théorème 4.3.6 (p. 70). On dit que c'est une \textit{catégorie $\pi$-$\mathfrak{U}$-paratopos}, si elle est $\pi$-accessible, i.e. si elle admet une petite sous-catégorie génératrice dont les objets soient $\pi$-accessible, ou ce qui revient au même, si % [sic] soient π-accessible -> soient π-accessibles, accord du pluriel
$\mathcal{M}$ est équivalente à une catégorie de la forme $\Ind'_{\pi}(\mathcal{M}_{\pi})$, où $\mathcal{M}_{\pi}$ est une petite catégorie (ou : une catégorie équivalente à une petite catégorie), et stable par $\varinjlim$ de cardinal $\leqslant \pi$. (Et on peut prendre pour $\mathcal{M}_{\pi}$ la sous-catégorie pleine de $\mathcal{M}$ formée des objets
\msp{page 107}
$\pi$-accessibles).

Dans la suite, nous allons, plus généralement, considérer une catégorie $C$ équivalente à une petite catégorie, stable par $\varinjlim$ de cardinal $\leqslant \pi_0$ (\footnote{Il n'y a pas lieu de le supposer tout de suite.}), et la catégorie
\[
\mathcal{M} \ \overset{\text{déf}}{=}\ \Ind'_{\pi_0}C \ \subset\ C^{\wedge}\ ,
\tag*{(4.5.1)}
\]
où $\pi_0$ est un cardinal infini donné. Je me propose d'introduire sur $\mathcal{M}$ une filtration cardinale canonique :

\textbf{Proposition 4.5.2} (\footnote{N'est prouvé (et n'est vrai) que si $\pi^{\pi_0} = \pi$, voir rectification 4.5.6, page 116.}). \textit{Soit $F \in \Ob\mathcal{M} \subset \Ob C^{\wedge}$, et soit $\pi$ un cardinal $\geqslant \pi_1 =$} $\sup(\pi_0, \card\Fl C_0)$\textit{, où $C_0$ est une petite sous-catégorie pleine de $C$ telle que $C_0 \overset{\sim}{\hookrightarrow} C$. Les conditions suivantes sont équivalentes :}
\begin{enumerate}[leftmargin=3em]
\item[\textit{a)}] $\card\Fl(C_0/F) \leqslant \pi$.
\item[\textit{a$'$)}] $\card\Ob\,(C_0/F) \leqslant \pi$.
\item[\textit{a$''$)}] \textit{Pour tout $x$ dans $C_0$,} $\card F(x) \leqslant \pi$.
\item[\textit{b)}] \textit{Il existe une petite catégorie filtrante $I$ (pas nécessairement ordonnée),} grande devant $\pi_0$\textit{, avec} $\card I \leqslant \pi$\textit{, et un système inductif $(x_i)_{i\in I}$ dans $C$, tel que l'on ait}
\[
F \ \simeq\ \varinjlim_{i\in I} x_i\ ,
\tag*{(4.5.2.1)}
\]
\textit{où la $\varinjlim$ est prise dans $\mathcal{M}$, ou ce qui revient au même (en vertu de 4.4.15, $I$ étant grande devant $\pi_0$) dans $C^{\wedge}$, i.e. $F$ est isomorphe au foncteur ind-représentable défini par l'ind-objet $\underline{x} = (x_i)_{i\in I}$ de $C$.}
\item[\textit{c)}] \textit{$F$ est $\pi$-accessible dans $\mathcal{M}$.}
\item[\textit{c$'$)}] \textit{$F$ est $\pi$-accessible dans $C^{\wedge}$.}
\end{enumerate}
\msp{page 108}

\textsc{Démonstration}. On a les implications $\pm$ tautologiques
\[
\begin{tikzcd}[row sep=small]
 & \mathrm{a}' \arrow[r, Rightarrow] & \mathrm{a}'' & \\
\mathrm{a} \arrow[ur, Rightarrow] \arrow[dr, Rightarrow] & & & \\
 & \mathrm{b} \arrow[r, Rightarrow] & \mathrm{c}' \arrow[r, Rightarrow] & \mathrm{c}\ ,
\end{tikzcd}
\tag*{$(*)$}
\]
\printedpage{30}
% head: A. Grothendieck : Dérivateurs, Ch. XVIII, Édition des Universités de Montpellier II et Paris VII
% foot: Version 17 décembre 2025
où a$'$ $\Longrightarrow$ a$''$ provient de l'inclusion canonique $F(x) = \Hom_{\mathcal{M}}(x,F) \hookrightarrow \Ob C/F$, a $\Longrightarrow$ b s'obtient en prenant $I = C_0/F$ et notant que $I$ est stable par limites inductives de cardinal $\leqslant \pi_0$ ($C$ l'étant), donc est grande devant $\pi_0$ (\footnote{On n'a pas besoin de cette stabilité, de toutes façons $C_0/F$ est grande devant $\pi_0$, en vertu de 4.4.3, p. 80.}). On a b $\Longrightarrow$ c$'$ en vertu du fait que la sous-catégorie strictement pleine $C^{\wedge}_{\pi}$ de $C^{\wedge}$ définie par
\[
\left\{
\begin{array}{lcl}
C^{\wedge}_{\pi} & = & \text{sous-catégorie strictement pleine de } C^{\wedge} \text{ formée des}\\
 & & F \in \mathcal{M} \text{ qui sont } \pi\text{-accessibles.}
\end{array}
\right.
\tag*{(4.5.2.2)}
\]
est stable par $\varinjlim$ de cardinal $\leqslant \pi$ (SGA 4 I 9.9). Enfin c$'$ $\Longrightarrow$ c provient du fait que le foncteur d'inclusion $\mathcal{M} \hookrightarrow C^{\wedge}$ est pleinement fidèle et $\pi_0$-accessible, i.e. commute aux limites inductives grandes devant $\pi_0$, et a fortiori à celles grandes devant $\pi$. Par les implications du diagramme ci-dessus, il suffit donc de prouver qu'on a
% marginal mark (handwritten squiggle) in the left margin of the edition beside the following display
\[
\mathrm{c} \overset{(\mbox{\footnotemark[33]})}{\Longrightarrow} a \qquad\quad \text{et} \qquad\quad \mathrm{a}'' \ \Longrightarrow\ \mathrm{a}\ .
\]
\footnotetext[33]{n'est prouvé (et n'est vrai) que si $\pi^{\pi_0} = \pi$.}
\setcounter{footnote}{33}

\msp{page 109}

\textit{Prouvons} a$''$ $\Longrightarrow$ a. Considérons l'application
\[
\Fl C/F \ \longrightarrow\ \Fl C
\tag*{$(*)$}
\]
induit par le foncteur étale $C' = C/F \longrightarrow C$, alors l'ensemble des flèches au-dessus d'une % [sic] induit -> induite, accord avec « l'application »
flèche $\alpha : x \longrightarrow y$ de $C$ est en correspondance biunivoque avec $\Ob C'_y = F(y)$, puisqu'une telle flèche $\alpha'$ est connue quand on connaît son but, qui est un objet quelconque de $F(y)$ ($C' \longrightarrow C$ étant étale, i.e. fibrant à fibres discrètes). Donc le fibre de $(*)$ en $\alpha$ est de % [sic] le fibre -> la fibre, genre
cardinal $= \card F(y) \leqslant \pi$, donc on a
\[
\card\Fl C/F \ \leqslant\ \pi \underbrace{\card\Fl C}_{\leqslant \pi_1} \ \leqslant\ \pi \underbrace{\ \pi_1\ }_{\leqslant \pi} \ \leqslant\ \pi^2 = \pi\ ,
\]
ce qui prouve a.

Prouvons enfin c $\Longrightarrow$ a. Pour ceci, nous allons utiliser la première inclusion $\mathcal{M}_{\pi} \subset \Filt^{\pi}\mathcal{M}$ de SGA 4 I 9.16, qui s'applique dès lors que nous prouvons le

% marginal mark (handwritten squiggle) in the left margin of the edition beside Corollaire 4.5.3
\textbf{Corollaire 4.5.3} (\footnote{Pas prouvé, la condition c) des filtrations cardinales peut-être pas satisfaite. Sans doute \textit{faux} en général.}). \textit{Sous les conditions de 4.5.2, soit, pour tout cardinal $\pi \geqslant \pi_1$,} $\Filt^{\pi}(\mathcal{M})$ \textit{la sous-catégorie strictement pleine de $\mathcal{M}$ formée des $F \in \Ob\mathcal{M}$ qui satisfont la condition a) de 4.5.2. Alors la famille des} $(\Filt^{\pi})_{\pi\geqslant\pi_1}$ \textit{est une}
\msp{page 110}
\textit{filtration cardinale de $\mathcal{M}$.}

La condition a) de loc. cit. 9.12 (les $\uFilt^{\pi}$ sont équivalentes à des petites catégories) est évidente. Prouvons b), i.e. la stabilité de $\uFilt^{\pi}$ par $\varinjlim$ filtrantes $F = \varinjlim F_i$ \textit{grandes devant} $\pi_0$, de cardinal $\leqslant \pi$. Comme $\mathcal{M} \hookrightarrow C^{\wedge}$ commute à ces limites, pour $x$ dans $C$
\[
F(x) \ =\ \varinjlim_{i} F_i(x)\ ,
\]
\printedpage{31}
% head: A. Grothendieck : Dérivateurs, Ch. XVIII, Édition des Universités de Montpellier II et Paris VII
% foot: Version 17 décembre 2025
d'où $\card F(x) \leqslant \card\Ob I \cdot \pi \leqslant \pi \cdot \pi = \pi$, ce qui (par l'équivalence a $\Longleftrightarrow$ a$''$) implique $F \in \Ob\Filt^{\pi}$. Prouvons la condition c) de loc. cit. Pour ceci, nous aurons besoin du fait, plus précis que le b), qu'on vient d'établir :

\textbf{Corollaire 4.5.4} (\footnote{Prouvé seulement si $\pi^{\pi_0} = \pi$\,! Sans doute \textit{faux} en général.}). \textit{Pour tout $\pi \geqslant \pi_1$,} $\Filt^{\pi}$ \textit{est stable par $\varinjlim$ (dans $\mathcal{M}$) de cardinal $\leqslant \pi$ (pas nécessairement filtrantes, a fortiori pas nécessairement grandes devant $\pi_0$).}

Il suffit, en vertu de ce qui est déjà prouvé de du lemme 4.4.14 (\footnote{Il faut une variante du lemme 4.4.14, pour les $\varinjlim$ de cardinal $\leqslant \pi'$, qui se prouve de la même façon.}), de montrer que $\Filt^{\pi}$ % [sic] de du lemme -> du lemme, mot répété
est stable par $\varinjlim$ \textit{de cardinal} $\leqslant \pi_0$. Si $\varphi : J \longrightarrow \Filt^{\pi}$ est un foncteur avec $\card J \leqslant \pi_0$, on le met sous forme indicielle
\[
\psi\ :\ I \times J \ \longrightarrow\ C\ ,
\]
grâce à 4.14.3.b, avec $I$ grande devant $\pi_0$.

\msp{page 111}

Je rappelle qu'on pouvait prendre
\[
I \ =\ \uHom(J,C)/\varphi\ ,
\]
où $\uHom(J,C) \hookrightarrow \uHom(J, \mathcal{M} = \Ind'_{\pi}C)$, quitte à ne pas exiger que $I$ soit ordonnée. On a envie de prouver que $\card\Fl I \leqslant \pi$. L'ennui, c'est qu'on arrive seulement à la majoration
\[
\card\Fl I \ \leqslant\ \pi^{\pi_0}\ .
\]
(Même si $J$ est discrète, on n'a pas mieux, puisqu'alors
\[
I \ \simeq\ \prod_{j\in J} C/\varphi(j)\ , \qquad \text{d'où} \quad \Fl I = \prod_{j\in J} \Fl C/\varphi(j)\ ,
\]
et si $\card\Fl C/\varphi(j) = \pi$ pour tout $j$, on trouve $\card\Fl I = \pi^{\pi_0}$ -- il n'y a pas moyen de faire mieux\,!)

Donc je n'arrive pas à prouver directement 4.5.4, ce qui m'aurait arrangé pour prouver la condition c) de loc. cit. Mais à vrai dire, on a fait ici le travail de la démonstration du théorème 4.3.6, partie b $\Longrightarrow$ c -- on y a déjà prouvé que la filtration qu'on vient de construire (par la condition a) de 4.5.2) est une filtration cardinale. La condition c) de SGA 4 I 9.12 est prouvée à la page 76. Mais je vais refaire
\msp{page 112}
la démonstration plus sérieusement.

Soit donc $F \in \Ob\mathcal{M}$, considérons
\[
J \ =\ C/F\ ,
\]
c'est une catégorie stable par $\varinjlim$ de cardinal $\leqslant \pi_0$, donc grande devant $\pi_0$, et on a
\[
F \ \simeq\ \varinjlim_{j\in J} x_j \qquad\quad (\,\varinjlim \ \text{dans } \mathcal{M})
\]
\printedpage{32}
% head: A. Grothendieck : Dérivateurs, Ch. XVIII, Édition des Universités de Montpellier II et Paris VII
% foot: Version 17 décembre 2025
en désignant par $j \longmapsto x_j$ le foncteur canonique $C/F \longrightarrow C$. Soit alors $I$ l'ensemble des sous-catégories $J'$ de $J$ qui sont \textit{de cardinal $\leqslant \pi$ et grandes devant $\pi_0$, $I$} ordonné par inclusion, et pour tout $J' \in I$, posons
\[
X_{J'} \ =\ \varinjlim_{j\in J'} x_j\ .
\]
Comme $J'$ est grande devant $\pi_0$, la limite inductive dans $\mathcal{M}$ est aussi une limite dans $C^{\wedge}$, d'où résulte aussitôt, comme $\card\Ob J' \leqslant \card\Fl J' \leqslant \pi$, que $X_{J'}$ est dans $\Filt^{\pi}$. Alors
\[
J' \ \longmapsto\ X_{J'} \ \in\ \Ob\Filt^{\pi}
\]
est un foncteur
\[
I \ \longrightarrow\ \Filt^{\pi}\ ,
\]
dont la limite inductive est celle des $x_j$, donc isomorphe à $F$. D'autre part, il est immédiat que $I$ est grand devant $\pi$, i.e. que si $(J'_{\alpha})_{\alpha\in A}$ est
\msp{page 113}
une famille d'éléments de $I$, celle-ci est majorée par un $J' \in I$. Ceci résulte du

% marginal mark (handwritten squiggle) in the left margin of the edition beside Lemme 4.5.5
\textbf{Lemme 4.5.5} (\footnote{N'est prouvé \textit{que} si $\pi^{\pi_0} = \pi$\,! Sans doute faux en général.}). \textit{Soient $\pi$, $\pi_0$, avec $\pi > \pi_0$, des cardinaux infinis, et soit $J$ une catégorie grande devant $\pi_0$. Alors pour toute partie $\Phi$ de} $\Fl J$\textit{, telle que} $\card\Phi \leqslant \pi$\textit{, il existe une sous-catégorie $J'$ de $J$} grande devant $\pi_0$\textit{, telle que} $\card J' \leqslant \pi$ (\footnote{C'est OK si on exige seulement $\card J' \leqslant \pi^{\pi_0}$.})\textit{, et} $\Fl J' \supset \Phi$.

Ce lemme étant admis, on a donc représenté un objet quelconque $F$ de $\mathcal{M}$ comme une $\varinjlim\limits_{i\in I} X_i$, avec $I$ ensemble ordonné grand devant $\pi$, les $X_i$ dans $\Filt^{\pi}$. Il reste à prouver que si $\pi' > \pi$ et $F \in \Filt^{\pi'}$, alors on peut prendre $I$ tel que $\card I \leqslant \pi'^{\pi}$. Mais $F \in \Filt^{\pi'}$ signifie que $J = C/F$ satisfait $\card\Fl J \leqslant \pi'$. Alors l'ensemble des parties de cardinal $\leqslant \pi$ de $\Fl J$ a un cardinal majoré par $\pi'^{\pi}$, a fortiori $\card I \leqslant \pi'^{\pi}$, ce qui achève de prouver que $(\Filt^{\pi})_{\pi\geqslant\pi_1}$ est une filtration cardinale sur $\mathcal{M}$.

\msp{page 114}

Pour prouver le lemme, on commence par saturer $\Phi$ par composition finie, on trouve une sous-catégorie $J'_0$ de $J$ telle que $\card\Fl J'_0 \leqslant \pi_0$ \ldots{} Mais je m'aperçois que pour prouver le lemme, j'ai encore besoin de $\pi^{\pi_0} = \pi$ -- une enveloppe de $J'$ grande devant $\pi_0$ va avoir un cardinal majoré (au mieux) par $\pi^{\pi_0}$.

Prenons p.ex.
\[
J \ =\ \mathfrak{P}(S) \qquad\quad \text{(ensemble ordonné)}
\]
avec $\card S = \pi > \pi_0$, donc $J$ est grand devant $\pi_0$. Considérons l'application
\[
S \ \longrightarrow\ J = \mathfrak{P}(S)\ , \qquad\quad s \longmapsto \{s\}\ ,
\]
et prenons pour $\Phi \subset \Fl J$ l'ensemble des $\mathrm{id}_{\{s\}}$. Donc on cherche une partie $J'$ de $J$, contenant $S$, et grande devant $\pi_0$ pour l'ordre induit par $J$. On voit tout de suite que la plus petite partie $J'$ faisant l'affaire est l'ensemble $\mathfrak{P}_{\pi_0}(S)$ des parties de $S$ de cardinal $\leqslant \pi_0$. Mais on trouve
\[
\card J' \cdot \pi_0! \ =\ \pi^{\pi_0}\ ,
\]
\printedpage{33}
% head: A. Grothendieck : Dérivateurs, Ch. XVIII, Édition des Universités de Montpellier II et Paris VII
% foot: Version 17 décembre 2025
\msp{page 115}
où
\[
\begin{aligned}
\pi_0! \ &=\ \card\Aut\,(E_0) \qquad\quad (E_0 \text{ ensemble de cardinal } \pi_0)\\
&=\ \pi_0^{\pi_0}\\
&=\ 2^{\pi_0}\ ,
\end{aligned}
\]
donc
\[
\card J' \cdot 2^{\pi_0} \ =\ \pi^{\pi_0}\ .
\tag*{$(*)$}
\]
D'autre part, il est immédiat que $\card J' \geqslant \card\mathfrak{P}(E_0) = 2^{\pi_0}$ (où $E_0$ est une partie de $S$ de cardinal $\pi_0$), donc le premier membre de $(*)$ est égal à $\card J'$, donc on trouve
\[
\card J' \ =\ \pi^{\pi_0}\ ,
\]
sorry\,!

Finalement, je me suis convaincu que la proposition 4.5.2, et les corollaires et lemmes précédents par lesquels j'ai essayé de l'établir, n'est vrai que pour ces cardinaux $\pi > \pi_0$ tels que
\[
\pi^{\pi_0} \ =\ \pi\ ,
\]
ce qui, moyennant l'hypothèse du continu généralisée, équivaut à
\[
\pi^{\pi_0} \ <\ 2^{\pi}\ (\ =\ \pi^{\pi}\ )
\]
(ou encore à $\pi^{\pi_0} \neq 2^{\pi}$ $(= \pi^{\pi})$). Je vais donner
\msp{page 116}
ici un énoncé rectifié.

\textbf{Proposition 4.5.6.}
\begin{enumerate}[leftmargin=3em]
\item[\circled{1\textsuperscript{o}}] \textit{Sous les hypothèses générales de 4.5.2 sur $C$, $\pi_0$, $\pi_1$, $\pi$, les conditions a, a$'$, a$''$, b, c$'$ de cette proposition sont équivalentes (moyennant $\pi \geqslant \pi_1$)} (\footnote{\textbf{N.B.} Pour cette équivalence, l'hypothèse que $C$ soit stable par $\varinjlim$ de cardinal $\leqslant \pi$ est inutile, et aussi pour que c$'$) implique c), et lui soit équivalente si $\pi = \pi^{\pi_0} \geqslant \pi_1$, cf. 4.7.1, p. 161.})\textit{, et impliquent la condition c ($F$ est $\pi$-accessible). Quand $\pi^{\pi_0} = \pi$ (notamment quand $\pi = 2^{c}$, avec $c \geqslant \pi_0$), alors toutes les conditions a à c$'$ de 4.5.2 sont équivalentes.}
\item[\circled{2\textsuperscript{o}}] \textit{Plus généralement et plus précisément, pour $\pi \geqslant \pi_1$ donné, considérons les conditions sur $\pi$ :}
\begin{enumerate}[leftmargin=2em]
\item[\textit{1)}] $\pi^{\pi_0} = \pi$.
\item[\textit{2)}] $\Filt^{\pi}(\mathcal{M})$ \textit{(définie dans 4.5.3) est stable par $\varinjlim$ de cardinal $\leqslant \pi$.}
\item[\textit{3)}] \textit{$\forall$ $F$ dans $\mathcal{M}$, on a $F \simeq \varinjlim\limits_{I} X_i$, avec $I$ ensemble ordonné grand devant $\pi$, les $X_i$ dans} $\Filt^{\pi}(\mathcal{M})$\textit{. De plus, si $\pi' > \pi$ et $F \in$} $\Filt^{\pi'}(\mathcal{M})$\textit{, on peut prendre $I$ de cardinal $\leqslant \pi'^{\pi}$.}
\item[\textit{4)}] \textit{Comme 3), sans la majoration de cardinaux, et en supposant seulement $I$ catégorie grande devant $\pi$ (pas nécessairement catégorie ordonné).} % [sic] catégorie ordonné -> catégorie ordonnée, accord
\item[\textit{5)}] \textit{La condition c) de 4.5.2 implique les autres conditions (elles sont équivalentes, par 1°), en d'autres termes}
\[
\mathcal{M}_{\pi} \ \subset\ \Filt^{\pi} \qquad\quad \textit{(donc } \mathcal{M}_{\pi} = \Filt^{\pi}\textit{, puisqu'on sait que } \Filt^{\pi} \subset \mathcal{M}_{\pi}\textit{)}\ ,
\]
\end{enumerate}
\end{enumerate}

% ---- from der-ch18_p036-040.tex
\setcounter{footnote}{39}

\printedpage{34}
% head: A. Grothendieck : Dérivateurs, Ch. XVIII, Édition des Universités de Montpellier II et Paris VII
% foot: Version 17 décembre 2025

\msp{page 117}

\begin{quote}
\textit{où $\mathcal{M}_\pi$ désigne la sous-catégorie strictement pleine de $\mathcal{M}$ formée des objets $\pi$-accessibles.}
\end{quote}

\begin{quote}
\textit{Ceci posé, on a les implications}
\[
\mathrm{1)} \ \Longrightarrow\ \mathrm{2)} \ \Longrightarrow\ \mathrm{3)} \ \Longrightarrow\ \mathrm{4)} \ \Longrightarrow\ \mathrm{5)} \ ,
\]
\textit{et dans le cas particulier où $\mathcal{M} = $} (Ens), \textit{$C$ = sous-catégorie pleine formée des ensembles de cardinal $\leqslant \pi_0$, on a aussi} 5) $\Longrightarrow$ 1), \textit{i.e. toutes ces conditions sont équivalentes.}
\end{quote}

\textsc{Démonstration.}

1) $\Longrightarrow$ 2). Cf. démonstration de 4.5.4.

2) $\Longrightarrow$ 3). La catégorie $J = \Filt^\pi/F$ est stable par $\varinjlim$ de cardinal $\leqslant \pi$, $\Filt^\pi$ l'étant, donc $J$ est grande devant $\pi$, d'ailleurs $F = \varinjlim\limits_{j\in J} X_j$, où $j \longmapsto X_j$ est le foncteur source $\Filt^\pi/F \longrightarrow \Filt^\pi$. D'où la représentation voulue de $F$, quitte à remplacer $J$ par un ensemble ordonné cofinal $I$, p.ex. en prenant pour $I$ l'ensemble des sous-catégories $J$ de $I$ telles que $\card \Fl J \leqslant \pi$ \texttt{[cette deuxième partie de la phrase “p.ex. ...” demande correction]}. La majoration $\card I \leqslant \pi'^{\pi}$, pour $F \in \Ob \Filt^{\pi'}$, se fait par AQT \texttt{[âne qui trotte]}.

L'implication 3) $\Longrightarrow$ 4) est tautologique.

4) $\Longrightarrow$ 5) (\footnote{comparer le lemme 4.5.14, page 133.}), car si $F \in \Ob \mathcal{M}_\pi$, on trouve par 4) que $F$ est facteur direct d'un objet de $\Filt^\pi$, et $\Filt^\pi$ étant stable par $\varinjlim$
\msp{page 118}
grandes devant $\pi_0$, cela implique que $\Filt^\pi$ est stable par facteurs directs, d'où $F \in \Ob \Filt^\pi$.

Enfin, si $\mathcal{M} = \Ens$, $C = \Ens(\pi_0)$. On note que
\[
\Filt^\pi \ = \ \{ \ X \in \Ob \Ens \mid \card\,(X)^{\pi_0} \ \leqslant \ \pi \ \} \ .
\]
D'autre part, on trouve facilement
\[
\begin{aligned}
(\Ens)_\pi \quad &\overset{\text{déf}}{=} \quad \text{sous-catégorie pleine de (Ens) formée} \\
&\phantom{\overset{\text{déf}}{=}} \quad \text{des objets $\pi$-accessibles} \\
&= \quad \{ X \in \Ob \Ens \mid \card X \leqslant \pi \} \ .
\end{aligned}
\]
Donc on a $(\Ens)_\pi \subset \Filt^\pi$ si et seulement si $\pi^{\pi_0} \leqslant \pi$, i.e. $\pi^{\pi_0} = \pi$, q.e.d. (\footnote{Cf. démonstration de la partie 1°) de 4.5.6, plus bas (bas de page et page suivante).}).

\textbf{Corollaire 4.5.7.} \textit{Pour tout cardinal $\pi \geqslant \pi_1$, on a}
\[
\Filt^\pi \ \subset \ \mathcal{M}_\pi \ \subset \ \Filt^{\pi^{\pi_0}} \ . \tag*{(4.5.7.1)}
\]

L'inclusion $\Filt^\pi \subset \mathcal{M}_\pi$ est l'implication a $\Longrightarrow$ c de la partie 1°) de 4.5.6. Pour prouver $\mathcal{M}_\pi \subset \Filt^{\pi^{\pi_0}}$, on pose $\pi' = \pi^{\pi_0}$ et on note que $\pi'^{\pi_0} = \pi^{\pi_0 \cdot \pi_0} = \pi^{\pi_0} = \pi'$. Donc on peut
\printedpage{35}
% head: A. Grothendieck : Dérivateurs, Ch. XVIII, Édition des Universités de Montpellier II et Paris VII
% foot: Version 17 décembre 2025
appliquer 1° à $\pi'$ au lieu de $\pi$, et on trouve que $\mathcal{M}_{\pi'} = \Filt^{\pi'}$. Comme $\mathcal{M}_\pi \subset \mathcal{M}_{\pi'}$, on en conclut $\mathcal{M}_\pi \subset \Filt^{\pi'}$, q.e.d.

J'ai oublié de démontrer le 1°) dans 4.5.6. Vu le diagramme d'implications $(\ast)$, p. 108, et l'implication $\mathrm{a}'' \Longrightarrow \mathrm{a}$ déjà prouvé (p. 109),
\msp{page 119}
il reste à prouver seulement
\[
\mathrm{c}' \ \Longrightarrow \ \mathrm{a}'' \ ,
\]
i.e. que $F$ $\pi$-accessible \textit{dans} $C^\wedge$ implique $\card F(x) \leqslant \pi \ \forall \ x \in C$. Ceci va résulter du résultat un peu plus général :

\textbf{Lemme 4.5.8.}
\begin{itemize}[leftmargin=3em]
\item[\textit{1°)}] \textit{Dans la catégorie $\mathcal{E} = $} (Ens), \textit{les objets $\pi$-accessibles ($\pi$ un cardinal infini) sont les ensembles $X$ tels que} $\card X \leqslant \pi$.
\item[\textit{2°)}] (\footnote{Cf. généralisation 4.14.18, p. 391 b) et c).}). \textit{Soit $C$ une catégorie équivalente à une petite catégorie, telle que $\forall \ x, y \in \Ob C$, on ait} $\card \Hom(x,y) \leqslant \pi$. \textit{Soit $\mathcal{E}$ une $\mathfrak{U}$-catégorie stable par $\varinjlim$ grandes devant $\pi$, et faiblement $\pi$-accessible (i.e. admettant une sous-catégorie génératrice formée d'objets $\pi$-accessibles), et stable par produits de cardinal $\leqslant \pi$. Alors pour qu'un objet $F$ de $\uHom(C,\mathcal{E})$ soit $\pi$-accessible, il} faut \textit{que $\forall \ x \in \Ob C$, $F(x)$ soit un objet $\pi$-accessible de $\mathcal{E}$, et cette condition est aussi suffisante si} $\card \operatorname{\text{réd}} \Fl(C) \leqslant \pi$ (\footnote{même \textit{sans} condition de stabilité de $\mathcal{E}$ par produits de cardinal $\leqslant \pi$.}).
\end{itemize}

\textsc{Démonstration.} 1°) $\card X \leqslant \pi \Longrightarrow X$ $\pi$-accessible dans $\mathcal{E} = (\Ens)$, car $X = \coprod_{x \in X} \{x\}$, et on applique SGA 4 I 9.9. Inversement, supposons $X$ $\pi$-accessible, et écrivons $X \simeq \varinjlim\limits_{i \in I} X_i$, où $I$ est l'ensemble des parties de $X$ qui sont de cardinal $\leqslant \pi$. Comme $I$ est grande devant $\pi$, l'application identique $X \xrightarrow{\ \mathrm{id}\ } X$ se factorise par un ensemble $X_i$, donc $X_i = X$, donc $\card X = \card X_i \leqslant \pi$, q.e.d.

\msp{page 120}

2°) Considérons, pour $x \in \Ob C$, le foncteur
\[
e \xrightarrow{\ \varepsilon_x\ } C \ ,
\]
d'où un foncteur
\[
\varepsilon_x^* \ : \ \uHom(C,\mathcal{E}) \longrightarrow \mathcal{E} \ , \qquad F \longmapsto F(x) \ ,
\]
admettant un adjoint à droite
\[
\varepsilon_{x\,*} \ : \ \mathcal{E} \ \longrightarrow \ \uHom(C,\mathcal{E}) \ ,
\]
donné par
\[
\varepsilon_{x\,*}(E)(y) \ = \ E^{\Hom(x,y)}
\]
(sous réserve d'existence des produits de cardinal $\leqslant \pi$). Ce foncteur $\varepsilon_{x\,*}$ est $\pi$-accessible en vertu de loc. cit. 9.8.

Ceci posé, soit $F$ un objet $\pi$-accessible de $\uHom(C,\mathcal{E})$, on voit que $\varepsilon_x^*(F)$ est $\pi$-accessible, $\varepsilon_x^*$ ayant un adjoint à droite $\pi$-accessible :
\[
E \ \longmapsto \ \Hom_{\mathcal{E}}(\underbrace{F(x)}_{\varepsilon_x^*(F)}, E) \ \simeq \ \Hom_{\uHom(C,\mathcal{E})}(F, \varepsilon_{x\,*}(E))
\]
\printedpage{36}
% head: A. Grothendieck : Dérivateurs, Ch. XVIII, Édition des Universités de Montpellier II et Paris VII
% foot: Version 17 décembre 2025
est composé du foncteur $\pi$-accessible $\varepsilon_{x\,*}$, et du foncteur covariant $\pi$-accessible représenté par $F$ dans $\uHom(C,\mathcal{E})$.

Inversement, si les $F(x)$ sont $\pi$-accessibles, prouvons que $F$ l'est, moyennant $\card \Fl(C) \leqslant \pi$. Se fait par AQT \texttt{[âne qui trotte]}, je ne vais pas l'expliciter ici.

\msp{page 121}

\textbf{Proposition 4.5.9.} \textit{Soit $C$ une catégorie équivalente à une petite catégorie, et soit $\mathcal{M} = $} $\Ind'_{\pi_0}(C)$, \textit{$\pi_0$ un cardinal infini donné. On suppose $C$ stable par $\varinjlim$ de cardinal $\leqslant \pi_0$. Pour tout cardinal $\pi \geqslant \pi_0$, soit $\mathcal{M}_\pi \supset C$ la sous-catégorie pleine de $\mathcal{M}$ formée des objets $\pi$-accessibles de $\mathcal{M}$. Soit d'autre part} $\pi_1 = \sup(\pi_0, \card \Fl C)$. \textit{Alors $(\mathcal{M}_\pi)_{\pi \geqslant \pi_1}$ est une filtration cardinale de $\mathcal{M}$} (\footnote{\textbf{N.B.} Toutes les conditions des filtrations cardinales sont vérifiées à l'exclusion, quand $\pi \geqslant \pi_0$, de celle de la majoration de cardinaux $\card I \leqslant \pi'^{\pi}$ condition c).})\textit{, et de plus, pour $\pi \geqslant \pi_0$, $\mathcal{M}_\pi$ est stable dans $\mathcal{M}$ par $\varinjlim$ de cardinal $\leqslant \pi$.}

\textsc{Démonstration.} La stabilité de $\mathcal{M}_\pi$ est un rappel, cf. SGA 4 I 9.9. Comme on a
\[
\mathcal{M}_\pi \ \subset \ \Filt^{\pi^{\pi_0}}
\]
pour $\pi \geqslant \pi_1$ (cf. 4.5.7), on voit donc que les $\mathcal{M}_\pi$ sont équivalentes à des petites catégories (condition a) des filtrations cardinales). La condition b) des filtrations cardinales est contenue dans la propriété de stabilité des $\mathcal{M}_\pi$ qu'on vient de rappeler. Prouvons enfin la condition c). Soit donc $\pi \geqslant \pi_0$, et soit $x$ dans $\mathcal{M}$, considérons la catégorie
\[
I \ = \ \mathcal{M}_\pi/x \ ,
\]
elle est stable par $\varinjlim$ de cardinal $\leqslant \pi$, d'après la stabilité similaire de $\mathcal{M}_\pi$, donc $I$ est grande devant $\pi$ et équivalente à une petite catégorie, d'autre part, on a
\msp{page 122}
\[
x \ = \ \varinjlim_{i \in I} x_i \ ,
\]
les $x_i$ dans $\mathcal{M}_\pi$ ( $\varinjlim$ dans $\mathcal{M}$ des foncteurs $I = \mathcal{M}_\pi/x \longrightarrow \mathcal{M}_\pi$). Quitte à remplacer $I$ par une catégorie ordonnée cofinale, cela prouve la partie qualitative de c. Supposons enfin $x \in \Ob \mathcal{M}_{\pi'}$, où $\pi' > \pi$, montrons que l'on a alors
\[
\card I_0, \ \card I' \ \leqslant \ \pi'^{\pi} \qquad \text{(\footnotemark, \footnotemark)} \ , \tag*{(4.5.9.1)}
\]
\addtocounter{footnote}{-1}%
\footnotetext{où $I_0$, $I'_0$ sont des sous-catégories réduites de $I$, $I'$ telles que $I_0 \hookrightarrow I$, $I'_0 \hookrightarrow I'$ soient des équivalences \texttt{[Grothendieck a biffé l'indice 0 de $I'$ dans 4.5.9.1]}.}%
\stepcounter{footnote}%
\footnotetext{On ne le prouve que si $\pi \geqslant \pi_1$.}%
ce qui achèvera de prouver que l'on a une filtration cardinale. Il nous faut une majoration des cardinaux.

\textbf{Lemme 4.5.10} (\footnote{Donner énonce commun avec 4.5.11, et prouver d'abord 4.5.11.}). \textit{Soit $\pi$ un cardinal $\geqslant \pi_1$. Alors on a} % [sic] énonce -> énoncé, missing accent (in footnote 47)
\[
\card \Fl((\Filt^\pi)_0) \ \leqslant \ 2^\pi \ . \tag*{(4.5.10.1)}
\]

On a en effet, en désignant par $C^\wedge(\pi)$ la sous-catégorie strictement pleine de $C^\wedge$ formée des $F : C^\circ \longrightarrow \Ens$ avec $\card F(x) \leqslant \pi$ pour tout $x$ dans $C$ (i.e. $F$ se factorise par $\Ens(\pi)$), et par $C^\wedge(\pi)_0$ une
\printedpage{37}
% head: A. Grothendieck : Dérivateurs, Ch. XVIII, Édition des Universités de Montpellier II et Paris VII
% foot: Version 17 décembre 2025
\msp{page 123}
sous-catégorie réduite exhaustive,
\[
\card C^\wedge(\pi)_0 \ \leqslant \ 2^\pi \ , \tag*{(4.5.10.2)}
\]
d'où résulte (4.5.10.1), puisque pour $\pi \geqslant \pi'$, $\Filt^\pi$ est une sous-catégorie pleine de $C^\wedge(\pi')$. Pour établir (4.5.10.2), on note que l'ensemble des classes à isomorphie près de $\Ens(\pi')$ s'identifie à l'ensemble des cardinaux $\leqslant \pi'$, lequel ensemble est de cardinal $\leqslant \pi'$. Remplaçant $\Ens(\pi)$ par $\Ens(\pi')_0$, $C$ par $C_0$, on doit majorer $\card \Fl(\uHom(C_0, \Ens(\pi')_0))$. Or $\Ob \uHom(C_0, \Ens(\pi')_0)$ est contenu dans l'ensemble des applications de $\Fl C_0$ dans $\Fl(\Ens(\pi')_0)$, le premier est majoré par $\pi_1$, le deuxième par $\pi^\pi \times \pi \times \pi = \pi^{2\pi} = \pi^\pi = 2^\pi$ (\textbf{N.B.} $\pi < 2^\pi$, d'où $\pi^\pi \leqslant (2^\pi)^\pi = 2^{\pi^2} = 2^\pi$), donc $\Ob \uHom$ majoré par $(2^\pi)^{\pi_1} = 2^{\pi\pi_1} = 2^\pi$. D'autre part, si $F$, $G$ sont dans $\uHom(C_0, \Ens(\pi_0))$, $\Hom(F,G)$ est majoré par $(\pi^\pi)^{\pi_1} = \pi^\pi = 2^\pi$, donc $\Fl \uHom$ est majoré par $2^\pi \times 2^\pi \times 2^\pi$, q.e.d.

\textbf{Corollaire 4.5.11.} \textit{Soit $F \in \Ob \Filt^\pi$, $F' \in \Ob \Filt^{\pi'}$, avec $\pi' \geqslant \pi \geqslant \pi_1$. Alors}
\[
\card \Hom(F, F') \ \leqslant \ \pi'^{\pi} \ . \tag*{(4.5.11.1)}
\]

\msp{page 124}

C'est essentiellement la même démonstration, légèrement plus simple. (On aurait dû commencer par là.)

\textbf{Corollaire 4.5.12.} \textit{Si $\pi' \geqslant \pi \geqslant \pi_1$, on a pour $x \in \Ob \mathcal{M}_{\pi'}$}
\[
\card\,(\mathcal{M}_\pi/x)_0 \ \leqslant \ \pi'^{(\pi^{\pi_0})} \ . \tag*{(4.5.12)}
\]

En effet,
\[
\card \Ob\,(\mathcal{M}_\pi/x)_0 \ \overset{\text{par 4.5.11}}{\leqslant} \ \underbrace{\underbrace{\card \Ob\,(\mathcal{M}_\pi)_0}_{\leqslant 2^{(\pi^{\pi_0})} \text{par 4.5.10}} \times (\pi'^{\pi_0})^{\pi^{\pi_0}}}_{= \pi'^{(\pi^{\pi_0})}} \ ,
\]
i.e.
\[
\card \Ob\,(\mathcal{M}_\pi/x)_0 \ \leqslant \ \pi'^{(\pi^{\pi_0})} \ ,
\]
d'où
\[
\card \Fl(\mathcal{M}_\pi/x)_0 \ \leqslant \ \pi'^{(\pi^{\pi_0})} \times \pi'^{(\pi^{\pi_0})} \times (\pi'^{\pi_0})^{(\pi^{\pi_0})} \ \leqslant \ \pi'^{(\pi^{\pi_0})} \ .
\]
Cela donne donc (revenant à 4.5.9.1)
\[
\card I_0 \ \leqslant \ \pi'^{(\pi^{\pi_0})} \ .
\]

D'autre part, $I'$ peut se définir en termes de $I_0$ comme l'ensemble des sous-catégories de $I_0$ qui sont de cardinal $\leqslant \pi$, donc leur nombre est majoré par $(\card I_0)^\pi \ \leqslant \ (\pi'^{\pi^{\pi_0}})^\pi = \pi'^{\pi\pi^{\pi_0}} = \pi'^{(\pi^{\pi_0})}$.

Mais ce n'est pas tout à fait la majoration exigée, puisque
\[
\pi'^{\pi} \leqslant \pi'^{(\pi^{\pi_0})} \qquad \text{(\footnotemark)},
\]
\footnotetext{\textbf{N.B.} On a égalité si on a $\pi' = 2^{c'}$ avec $c' \geqslant \pi^{\pi_0}$, soit $\pi = 2^c$ avec $c \geqslant \pi_0$.}%
et qu'il est possible, a priori, que l'inégalité
\printedpage{38}
% head: A. Grothendieck : Dérivateurs, Ch. XVIII, Édition des Universités de Montpellier II et Paris VII
% foot: Version 17 décembre 2025
\msp{page 125}
soit stricte ! Apparemment, les ``majorations''
\[
\mathcal{M}_\pi \ \subset \ \Filt^{\pi^{\pi_0}} \ , \qquad \mathcal{M}_{\pi'} \ \subset \ \Filt^{\pi'^{\pi_0}}
\]
sont trop grossières. Aussi il s'agit de mieux appréhender les catégories $\mathcal{M}_\pi$, pour $\pi \geqslant \pi_0$ (\footnote{Mais on peut utiliser 4.5.10, 4.5.11 pour conclure que pour $x \in \Ob \mathcal{M}_{\pi'} \subset \Ob \Filt^{\pi'^{\pi_0}}$, on a $\card \Fl(C_0/x) \leqslant (\pi'^{\pi_0})^{\pi_0} = \pi'^{\pi_0}$, on aura donc $x = \varinjlim\limits_{J} x_j$ ($x_j \in \Ob C$, $\varinjlim$ filtrante, $\card \Fl J \leqslant \pi'^{\pi_0}$), et en introduisant l'ensemble $I$ des sous-catégories $J' \subset J$ qui sont de cardinal $\leqslant \pi$, et posant $x_{J'} = \varinjlim\limits_{j \in J'} x_j$, on trouve $x = \varinjlim\limits_{J' \in I} x_{J'}$, où cette fois $I$ est un \textit{ensemble ordonné} grand devant $\pi$, les $x_{J'} \in \mathcal{M}_\pi$, et $\card I \leqslant (\card J)^\pi \leqslant (\pi'^{\pi_0})^\pi = \pi'^{\pi\pi_0} = \pi'^{\pi}$, c'est la bonne majoration !}). Je vais hasarder une

\textbf{Proposition 4.5.13.} \textit{Les hypothèses et notations sont celles de 4.5.9. Soit $\pi$ un cardinal $\geqslant \pi_0$, et soit $x \in \Ob \mathcal{M}$. Considérons les conditions :}
\begin{itemize}[leftmargin=3em]
\item[\textit{a})] \textit{Il existe une catégorie $I$, avec} $\card \Fl I \leqslant \pi$, \textit{et un foncteur $i \longmapsto x_i : I \longrightarrow C$, tel que $x$ soit isomorphe à $\varinjlim x_i$.}
\item[\textit{a}$'$)] \textit{Comme a), avec $I$ filtrante.}
\item[\textit{a}$''$)] \textit{Comme a), avec $I$ ordonnée filtrante.}
\item[\textit{b})] \textit{Il existe un foncteur} cofinal \textit{$I \longrightarrow C/x$, avec} $\card \Fl I \leqslant \pi$.
\item[\textit{b}$'$)] \textit{La catégorie filtrante $C/x$ est ``essentiellement de cardinal $\leqslant \pi$'', i.e. il existe une catégorie ordonnée} \texttt{[non, cf. b$''$)]} filtrante $I$, $\card \Fl I \leqslant \pi$, \textit{et un foncteur cofinal $I \longrightarrow C/x$.}
\item[\textit{b}$''$)] \textit{Comme b$'$), avec de plus $I$ ordonnée.}
\item[\textit{c})] \textit{Il existe une famille de flèches $(u_\alpha)_{\alpha \in A}$, $x_\alpha \longrightarrow x$, épimorphique stricte, indexée par un ensemble $A$ tel que} $\card A \leqslant \pi$.
\item[\textit{d})] \textit{On a $x \in \Ob \mathcal{M}_\pi$, i.e. $x$ est $\pi$-accessible dans $\mathcal{M}$.}
\end{itemize}

\msp{page 126}

\textit{Ceci posé :}
\begin{itemize}[leftmargin=3em]
\item[\circled{1\textsuperscript{\textit{o}}}] \textit{Les six premières conditions $a$, $a'$, $a''$, $b$, $b'$, $b''$ sont équivalentes, soit $P(\pi)$, et impliquent les deux dernières,}
\[
\begin{tikzcd}[column sep=small]
 & P(\pi) \arrow[dl, Rightarrow] \arrow[dr, Rightarrow] & \\
c & & d \ .
\end{tikzcd}
\]
\textit{Si $\pi \geqslant \pi_1$, alors $P(\pi)$ est équivalente à $c$.}
\item[\circled{2\textsuperscript{\textit{o}}}] \textit{Soit $\Filt'^{\pi} \subset \mathcal{M}$ la sous-catégorie strictement pleine de $\mathcal{M}$ formée des objets qui satisfont la condition $P(\pi)$ ci-dessus, et $\Filt^\pi$ la sous-catégorie strictement pleine définie dans 4.5.3 (p. 109). On a alors}
\[
\Filt^\pi \ \subset \ \Filt'^{\pi} \ \subset \ \mathcal{M}_\pi \ , \tag*{(4.5.13.1)}
\]
\textit{et pour qu'on ait égalité des trois membres, il suffit qu'on ait $\pi^{\pi_0} = \pi$. (Cf. prop. 4.5.6, page 116). De plus, $\Filt'^{\pi}$ est stable dans $\mathcal{M}$ par} sommes \textit{de cardinal $\leqslant \pi$.}
\item[\circled{3\textsuperscript{\textit{o}}}] \textit{Les conditions suivantes sur $C$, $\pi_0$, $\pi$ sont équivalentes.}
\end{itemize}

% ---- from der-ch18_p041-045.tex
\setcounter{footnote}{49}

\printedpage{39}
% head: A. Grothendieck : Dérivateurs, Ch. XVIII, Édition des Universités de Montpellier II et Paris VII

\begin{itemize}
\item[\textit{1)}] $\Filt'^{\pi}$ \textit{est stable dans $\mathcal{M}$ par} $\varinjlim$ \textit{de cardinal $\leqslant \pi$.}
\item[\textit{1$'$)}] $\Filt'^{\pi}$ \textit{est stable par conoyaux de double-flèches.}
\item[\textit{2)}] \textit{Tout objet $x$ de $\mathcal{M}$ est de la forme $\varinjlim\limits_{I} x_i$, où $x_i$ dans} $\Filt'^{\pi}$\textit{, $I$ grande devant $\pi$.}
\item[\textit{3)}] \textit{La condition d) de 1°) implique $P(\pi)$ (donc lui est équivalente), i.e. on a}
\[ \mathcal{M}_{\pi} \ = \ \Filt'^{\pi}. \]
\end{itemize}

\msp{page 127}
\textit{Ces conditions sont satisfaites dans les trois cas suivants :}
\begin{itemize}
\item[\textit{A)}] \textit{$\pi^{\pi_0} = \pi$ (pour mémoire, c'est contenu dans 2°).}
\item[\textit{B)}] \textit{Les familles épimorphiques strictes dans $\mathcal{M}$ sont épimorphiques strictes universelles.}
\item[\textit{C)}] \textit{Les limites inductives filtrantes dans $\mathcal{M}$ commutent au changement de base (ce qui est le cas notamment si les} $\varinjlim$ \textit{filtrantes sont exactes à gauche)} (\footnote{Je suppose que cette condition revient à la même condition sur $C$, pour les $\varinjlim$ filtrantes de cardinal $\leqslant \pi$.}).
\end{itemize}

\textbf{Remarque.} La filtration de $\mathcal{M}$ pour les $\Filt'^{\pi}$ n'est autre que celle envisagée dans SGA 4 I 9.13, où on fait de plus l'hypothèse B) sur $\mathcal{M}$, qui après-coup n'était plus tellement de mon goût. Toujours est-il que d'après (4.5.13.1), cette filtration se rapproche plus de la ``bonne'' filtration par les $\mathcal{M}_{\pi}$, que celle que j'avais eu encore naguère tendance à lui préférer, celle par les $\Filt^{\pi}$. Elle a d'ailleurs sur cette dernière l'avantage, dans le cas envisagé B), d'être une filtration cardinale (ce qui n'est pas nécessairement le cas pour $\Filt^{\pi}$, du moins s'il existe des cardinaux infinis $\pi \geqslant \pi_0$ avec $\pi^{\pi_0} \neq \pi$). Je soupçonne maintenant qu'en général (et s'il existe des cardinaux infinis $\pi > \pi_0$ tels que $\pi^{\pi_0} \neq \pi$), en dehors des conditions A) B) C), la filtration par les $(\Filt'^{\pi})$ n'est \textit{pas} une
\msp{page 128}
filtration cardinale. En tous cas, 4.5.6 et la présente proposition, 3°, montrent que si soit $(\Filt^{\pi})_{\pi\geqslant\pi_2}$, soit $(\Filt'^{\pi})_{\pi\geqslant\pi_2}$ est une filtration cardinale ($\pi_2$ étant un cardinal donné $\geqslant \pi_1$), alors elle coïncide avec la filtration par les $\mathcal{M}_{\pi}$. Donc apparemment, dans tous les cas où je sais construire une filtration cardinale, sur un paratopos du moins (i.e. sur une catégorie $\mathcal{M}$ stable par petites $\varinjlim$ , et dont tous les objets sont accessibles -- l'existence d'une filtration cardinale impliquant de plus celle d'une petite sous-catégorie strictement génératrice \dots), celle-ci n'est autre que celle par les $\mathcal{M}_{\pi}$ -- chose dont je ne m'étais par aperçu en rédigeant SGA 4 I 9. % [sic] m'étais par aperçu -> m'étais pas aperçu, "par" for the negation "pas"

\textsc{Démonstration de 4.5.13.}

\textit{Prouvons 1°.} On a les implications tautologiques
\[
\begin{tikzcd}[row sep=large]
\mathrm{b}'' \arrow[r, Rightarrow] \arrow[d, Rightarrow] & \mathrm{b}' \arrow[r, Rightarrow] \arrow[d, Rightarrow] & \mathrm{b} \arrow[d, Rightarrow] \\
\mathrm{a}'' \arrow[r, Rightarrow] & \mathrm{a}' \arrow[r, Rightarrow] & \mathrm{a} \text{ ,}
\end{tikzcd}
\qquad\qquad \text{(\footnotemark) ,}
\]
\footnotetext{compte tenu que $C$ est strictement génératrice dans $\mathcal{M}$, donc $x \simeq \varinjlim\limits_{i\in C/x} x_i$.}
et pour prouver l'équivalence de ces six conditions, il suffit de prouver
\[ \mathrm{a} \ \Longrightarrow \ \mathrm{b}'' \ . \]
% foot: Version 17 décembre 2025

\printedpage{40}
% head: A. Grothendieck : Dérivateurs, Ch. XVIII, Édition des Universités de Montpellier II et Paris VII

\msp{page 129}
Si on a
\[ x \ \simeq \ \varinjlim_{I} x_i \ , \qquad\quad \card \Fl I \ \leqslant \ \pi \ , \]
soit $I'$ l'ensemble des \textit{sous-diagrammes} $J \subset I$ finis de $I$, ordonné par inclusion, donc $I'$ est un ensemble ordonné filtrant, et du fait que $C$ est stable par $\varinjlim$ finies, on trouve un foncteur
\[ J \longmapsto x_J \ = \ \varinjlim_{i\in J} x_i \ : \ I' \longrightarrow C \ , \]
et on aura
\[ x \ \simeq \ \varinjlim_{J\in I'} x_J \ . \]
D'autre part, on a
\[
\begin{aligned}
\card I' \ &= \ \card \mathfrak{P}_f( \ \underbrace{\Fl(I)}_{\text{de card $\leqslant \pi$}} \ ) \\
&\leqslant \ \pi + \pi^2 + \ \cdots \ + \pi^n + \ \dots \\
&\leqslant \ \pi\aleph_0 \ = \ \pi \ ,
\end{aligned}
\]
q.e.d. Pour établir 1°, il reste à établir
\[
\begin{tikzcd}[column sep=small]
& P(\pi) \arrow[dl, Rightarrow] \arrow[dr, Rightarrow] & \\
\mathrm{c} & & \mathrm{d} \text{ ,}
\end{tikzcd}
\]
et l'implication inverse $\mathrm{c} \Longrightarrow P(\pi)$ quand $\pi \geqslant \pi_1$. Mais $P(\pi) \Longrightarrow \mathrm{c}$ est une tautologie, et $P(\pi) \Longrightarrow \mathrm{d}$ puisque $\mathcal{M}_{\pi}$ est stable par $\varinjlim$ de cardinal $\leqslant \pi$. Reste à prouver que $\mathrm{c} \Longrightarrow P(\pi)$ si $\pi \geqslant \pi_1$.

\msp{page 130}
Supposons donc qu'on ait une famille épimorphique stricte
\[ x_{\alpha} \longrightarrow x \qquad\quad (\alpha \in A,\ \card A \leqslant \pi) \ , \]
cela signifie donc que $x$ est limite inductive d'un système inductif indexé par la catégorie $B$ ayant comme ensemble d'objets $A \amalg (A\times A)$, les flèches non identiques étant représentées par les flèches canoniques
\[
\begin{tikzcd}[column sep=small]
& x_{\alpha} \times_x x_{\beta} \arrow[dl, "p_{\alpha,\beta}"'] \arrow[dr, "q_{\alpha,\beta}"] & \\
x_{\alpha} & & x_{\beta} \text{ .}
\end{tikzcd}
\]

(\textbf{N.B.} $B$ n'est par une catégorie ordonnée, car pour l'objet $(\alpha,\alpha)$, représenté par $x_{\alpha}\times_x x_{\alpha}$, % [sic] n'est par une -> n'est pas une, "par" for the negation "pas"
il y a \textit{deux} flèches dans l'objet $\alpha$, représenté par $x_{\alpha}$. Mais on a évidemment $\card \Fl(B) \leqslant \pi$.) Si on savait que $C$ est stable par produits binaires et par passage à des sous-objets (ici les $x_{\alpha} \times_x x_{\beta} \ \subset \ x_{\alpha} \times x_{\beta}$), on aurait prouvé la condition a). À défaut, utilisons encore le fait que $C$ est génératrice par épimorphismes (même stricts) et choisissons pour tout couple $(\alpha,\beta)$ un ensemble épimorphique de flèches :
\[ \underbrace{x_{\alpha,\beta,\gamma}}_{\in \Ob C_0} \xrightarrow{u_{\alpha,\beta,\gamma}} x_{\alpha} \times_x x_{\beta} \ , \]
% foot: Version 17 décembre 2025

\printedpage{41}
% head: A. Grothendieck : Dérivateurs, Ch. XVIII, Édition des Universités de Montpellier II et Paris VII
où $C_0$ est une sous-catégorie réduite exhaustive de $C$. La flèche $u_{\alpha,\beta,\gamma}$ est connue quand on connaît les composés
\msp{page 131}
\[
\begin{tikzcd}
& x_{\alpha} \\
x_{\alpha,\beta,\gamma} \arrow[ur, "v_{\alpha,\beta,\gamma}"] \arrow[dr, "w_{\alpha,\beta,\gamma}"'] & \\
& x_{\beta} \text{ ,}
\end{tikzcd}
\]
donc pour $x_{\alpha,\beta,\gamma} \in \Ob C_0$ fixé, les choix possibles pour $u_{\alpha,\beta,\gamma}$ sont de cardinal majoré par $\pi_1\pi_1 = \pi_1$, de plus on a $\card \Ob C_0 \leqslant \pi_1$, donc le nombre total des flèches $v_{\alpha,\beta,\gamma}$, $w_{\alpha,\beta,\gamma}$ (pour $\alpha$, $\beta$ fixés) est majoré par $\pi_1 \cdot \pi_1 = \pi_1$. Donc le nombre total des flèches $u_{\alpha,\beta,\gamma}$, pour $\alpha$, $\beta$, $\gamma$ variables, est majoré par $\pi \cdot \pi_1 = \pi$ (si $\pi \geqslant \pi_1$). Donc $x$ est représenté par une catégorie $I$ ayant comme objets l'ensemble somme de $A$ et de l'ensemble des indices triples $\alpha,\beta,\gamma$, et pour chaque tel indice, exactement deux flèches $v_{\alpha,\beta,\gamma}$, $w_{\alpha,\beta,\gamma}$ non identiques de source cet objet, et allant respectivement dans $x_{\alpha}$, $x_{\beta}$. On a donc $\card \Fl I \leqslant \pi$, et on a un foncteur $I \longrightarrow C$, $\alpha \longmapsto x_{\alpha}$, $(\alpha,\beta,\gamma) \longmapsto x_{\alpha,\beta,\gamma}$, tel que $x \ = \ \varinjlim\limits_{i\in I} x_i$. Cela achève la démonstration de 1°).

\msp{page 132}
\textit{Prouvons 2°.} La première inclusion (4.5.13.1) est tautologique, la deuxième ne fait qu'exprimer $P(\pi) \Longrightarrow \mathrm{d}$. Le fait que $\pi^{\pi_0} = \pi$ implique l'égalité des membres extrêmes est un rappel de 4.5.6. Reste à prouver que $\Filt'^{\pi}$ est stable dans $\mathcal{M}$ par sommes de cardinal $\leqslant \pi$. Mais si $(x_{\alpha})_{\alpha\in A}$ est une famille d'objets de $\Filt'^{\pi}$, avec $\card A \leqslant \pi$, et si $x = \coprod_{\alpha\in A} x_{\alpha}$, prenant pour tout $x_{\alpha}$ un foncteur $\varphi_{\alpha} : I_{\alpha} \longrightarrow C$, avec $\card \Fl I_{\alpha} \leqslant \pi$, tel que $x_{\alpha}$ soit la $\varinjlim$ dans $\mathcal{M}$ de ce foncteur, on trouve que $x$ est $\varinjlim$ du foncteur $\varphi : I \longrightarrow C$ défini par les $\varphi_{\alpha}$, où $I = \coprod_{\alpha} I_{\alpha}$ (somme dans Cat). Il est évident que $\card \Fl(I) \leqslant \pi\cdot\pi = \pi$, ce qui achève de prouver 2°.

\textit{Prouvons 3°.} On a les implications tautologiques $1 \Longrightarrow 1'$, et $1' \Longrightarrow 1$, puisque $\Filt'^{\pi}$ est stable par sommes de cardinal $\leqslant \pi$. Donc $1 \Longleftrightarrow 1'$, et implique que pour tout $x$ dans $\mathcal{M}$,
\msp{page 133}
$I \ = \ \Filt'^{\pi}/x$ est stable par $\varinjlim$ de cardinal $\leqslant \pi$, donc est grand devant $\pi$. Comme $x \simeq \varinjlim\limits_{i\in I} x_i$, cela prouve que $1 \Longrightarrow 2$. L'implication $2 \Longrightarrow 3$ est standard, et résulte de $\Filt'^{\pi} \subset \mathcal{M}_{\pi}$, je vais expliciter dans un

\textbf{Lemme 4.5.14.} \textit{Soient $\mathcal{M}$ une catégorie stable par petites} $\varinjlim$ \textit{, $\pi$ un cardinal infini, $\mathcal{M}_0$ une sous-catégorie strictement pleine de $\mathcal{M}$. Considérons les conditions}
\begin{itemize}
\item[\textit{a)}] \textit{$\mathcal{M}_0$ est une sous-catégorie strictement génératrice, et stable dans $\mathcal{M}$ par} $\varinjlim$ \textit{de cardinal $\leqslant \pi$.}
\item[\textit{b)}] \textit{Pour tout objet de $\mathcal{M}$, on a un isomorphisme}
\[ x \ \simeq \ \varinjlim_{i\in I} x_i \ , \]
\textit{les $x_i$ dans $\mathcal{M}_0$, $I$ grande devant $\pi$.}
% foot: Version 17 décembre 2025

\printedpage{42}
% head: A. Grothendieck : Dérivateurs, Ch. XVIII, Édition des Universités de Montpellier II et Paris VII
\item[\textit{b$'$)}] \textit{Comme b), avec $I$ catégorie} ordonnée.
\item[\textit{c)}] \textit{$\mathcal{M}_{\pi} \subset \mathcal{M}_0$ (où $\mathcal{M}_{\pi}$ est la sous-catégorie strictement pleine de $\mathcal{M}$ formée des objets $\pi$-accessibles).}
\end{itemize}
\textit{On a donc les implications}
\[ a \ \Longrightarrow \ b \ \Longleftrightarrow \ b' \ \Longrightarrow \ c \ , \]
\textit{et si on suppose $\mathcal{M}_0 \subset \mathcal{M}_{\pi}$, et $\mathcal{M}_{\pi}$ sous-catégorie strictement génératrice de $\mathcal{M}$, alors les quatre conditions $a$, $b$, $b'$, $c$ sont équivalentes, et équivalent à $\mathcal{M}_0 = \mathcal{M}_{\pi}$} (\footnote{cf. démonstration de proposition 4.5.6, $4 \Longrightarrow 5$, page 117.}).

\msp{page 134}
Enfin (revenant à la démonstration de 4.5.13, 3°), l'implication $3) \Longrightarrow 1)$ provient du fait que $\mathcal{M}_{\pi}$ est stable par $\varinjlim$ de cardinal $\leqslant \pi$.

Reste à prouver que ces conditions, i.e. la stabilité de $\Filt'^{\pi}$ par conoyaux de double-flèches, est satisfaite dans les cas A) B) C), le cas A) étant d'ailleurs déjà connu. Le cas B) résulte immédiatement de la transitivité des familles épimorphiques strictes universelles, et c'est d'ailleurs aussi contenu explicitement (stabilité de $\Filt'^{\pi}$ par $\varinjlim$ de cardinal $\leqslant \pi$) dans SGA 4 I 9.13 (sinon dans l'énoncé, du moins c'est signalé dans la démonstration). En rédigeant SGA 4 I 9, je savais donc qu'on a $\mathcal{M}_{\pi} \subset \Filt'^{\pi}$ pour tout $\pi$, mais l'inclusion, plus évidente pourtant dans le cas d'espèce, $\Filt'^{\pi} \subset \mathcal{M}_{\pi}$ (quand on suppose les objets de $C$ $\pi_1$-accessibles) m'avait apparemment échappé\,?

Reste à traiter le cas C). $\Big($J'ai l'impression qu'il est un peu plus général que le cas B), i.e. que la condition B) sur $\mathcal{M}$ implique C), mais non inversement. Ainsi, la condition C) est vérifiée dès que $\mathcal{M}$ correspond à une espèce de structure algébrique définie en
\msp{page 135}
termes de $\varprojlim$ \textit{finies}, i.e. quand $\mathcal{M}$ est de la forme $\Ind(C)$, $C$ stable par $\varinjlim$ \textit{finies}, alors que je doute que la condition B) soit toujours satisfaite pour ces catégories. Par exemple dans la catégorie Cat (qui est un paratopos de cette espèce), les $\varinjlim$ filtrantes sont exactes à gauche, mais je ne crois pas que les familles épimorphiques stricts soient universelles. % [sic] épimorphiques stricts -> épimorphiques strictes, agreement with "familles"
Prenons p.ex. deux homomorphismes de groupes $G_0 \xrightarrow{u_0} G$, $G_1 \xrightarrow{u_1} G$ tels que $G_0 * G_1 \longrightarrow G$ soit surjectif, alors cette famille à deux éléments est épimorphique stricte dans Gr et dans Cat (\footnote{!, c'est \textit{faux} en général, mais vrai dans le cas particulier plus bas.}), mais n'est pas forcément épimorphique stricte universelle, même si p.ex. $G_0 = G_1 = H$, $G = H \times H$, et $u_0, u_1$ sont des injections $H \rightrightarrows H \times H$, données par $u_0(h) = (h,1)$, et $u_1(h) = (1,h)$. Si je prends l'application diagonale $H \longrightarrow H \times H$, les images inverses de $u_0$, $u_1$ sont le même objet de $\mathcal{M}/H$, savoir $e \longrightarrow H$, et cette famille n'est épimorphique stricte que si $H$ est le groupe unité.$\Big)$

Le cas C) va résulter du résultat plus précis :

\textbf{Corollaire 4.5.15.} \textit{Sous les conditions générales de 4.5.13, supposons que dans $\mathcal{M}$ les} $\varinjlim$
\msp{page 136}
filtrantes \textit{commutent au changement de base (p.ex. les} $\varinjlim$ \textit{filtrantes sont exactes). Alors pour tout $\pi \geqslant \pi_1$,} $\Filt'^{\pi}$ \textit{est formée des $x \in \Ob \mathcal{M}$ satisfaisant la condition suivante.}
% foot: Version 17 décembre 2025

\printedpage{43}
% head: A. Grothendieck : Dérivateurs, Ch. XVIII, Édition des Universités de Montpellier II et Paris VII
\begin{itemize}
\item[\textit{c$'$)}] \textit{Il existe une famille de flèches $(u_{\alpha})_{\alpha\in A}$, $u_{\alpha} : x_{\alpha} \longrightarrow x$, dans $\mathcal{M}$, qui est} épimorphique stricte universelle \textit{(i.e. couvrante pour la topologie canonique de $\mathcal{M}$), avec $x_{\alpha} \in \Ob C$ $\forall\ \alpha \in A$, et} $\card A \leqslant \pi$.
\end{itemize}
\textsc{Démonstration} : il faut prouver que cette condition équivaut à la condition c) de 4.5.13 1°). Comme on a $\mathrm{c}' \Longrightarrow \mathrm{c}$ trivialement, il reste à montrer que $\mathrm{c} \Longrightarrow \mathrm{c}'$. Mais on a vu dans 4.5.13 1°) que c implique a$''$, i.e. qu'on a un isomorphisme
\[ x \ = \ \varinjlim_{i\in I} x_i \ , \]
les $x_i$ dans $C$, et $I$ ensemble \textit{ordonné filtrant}, avec $\card I \leqslant \pi$. Il reste à voir que si $x$ dans $\mathcal{M}$ est une limite inductive \textit{filtrante} d'objets $x_i$, alors l'ensemble des flèches $x_i \longrightarrow x$ est épimorphique stricte \textit{universelle.} Or cela est immédiat, car pour $x' \longrightarrow x$, on a $x' = \varinjlim\limits_{i} x'_i$ (par hypothèse C), où $x'_i = x' \times_x x_i$, donc la famille $x'_i \longrightarrow x'$ est épimorphique stricte, q.e.d.

Prouvons maintenant que dans le cas C, la condition de 3°, 1 est satisfaite, i.e. que $\Filt'^{\pi}$ est stable par $\varinjlim$ de cardinal $\leqslant \pi$. Cela
\msp{page 137}
résulte aussitôt du

\textbf{Corollaire 4.5.16.} (de 4.5.15). \textit{Sous les conditions de 4.5.15 (i.e. la condition C) de 4.5.13, 3°), soit $x$ un objet de $\mathcal{M}$, et supposons qu'il existe une famille} épimorphique \underline{universelle} $(x_{\alpha} \longrightarrow x)_{\alpha\in A}$ \textit{de flèches de $\mathcal{M}$, avec $x_{\alpha} \in$} $\Ob \Filt'^{\pi}$ \textit{pour tout $\alpha \in A$, et} $\card A \leqslant \pi$. \textit{Alors $x \in$} $\Ob \Filt'^{\pi}$.

En effet, l'hypothèse sur $x_{\alpha}$ équivaut, en vertu de 4.5.15, à l'existence d'une famille épimorphique stricte universelle de flèches $(x_{\alpha\beta} \longrightarrow x_{\alpha})_{\beta\in B_{\alpha}}$, avec $\card B_{\alpha} \leqslant \pi$ et les $x_{\alpha\beta}$ dans $C$. Soit $B = \coprod_{\alpha\in A} B_{\alpha}$, on a alors $\card B \leqslant \pi^2 = \pi$, et d'autre part, la famille des flèches composées $x_{\alpha\beta} \longrightarrow x_{\alpha} \longrightarrow x$ indexée par $B$ est épimorphique stricte universelle (par transitivité de ce type de familles), donc $x \in \Ob \Filt'^{\pi}$, q.e.d.

Avec tout ça, on n'a toujours pas donné de caractérisation de $\mathcal{M}_{\pi}$ lui-même, qui est meilleure que $\Filt^{\pi}$ et $\Filt'^{\pi}$ du point de vue catégorique, car seuls les $\mathcal{M}_{\pi}$ définissent une filtration cardinale, et de plus ils ont les propriétés de stabilité désirables dans ce
\msp{page 138}
contexte, pour les $\varinjlim$ de cardinal $\leqslant \pi$. Le seul critère non tautologique que j'arrive à dégager dans le cas général (pour tout $\pi \geqslant \pi_0$, et sans hypothèse d'exactitude supplémentaire sur $\mathcal{M}$) est le suivant, beaucoup moins ``explicite'' ou ``calculatoire'' que les critères très concrets pour $\Filt^{\pi}$ et pour $\Filt'^{\pi}$ :

\textbf{Proposition 4.5.17.} \textit{Soit $\pi$ un cardinal $\geqslant \pi_0$. Alors $\mathcal{M}_{\pi}$ est égal à la plus petite sous-catégorie strictement pleine de $\mathcal{M}$, contenant $C$, et stable par} $\varinjlim$ \textit{de cardinal $\leqslant \pi$.} (Il est important de noter qu'il \textit{ne suffit pas} de dire : stable par $\varinjlim\limits_{I}$ , avec \textit{$I$ grande devant} $\pi_0$, et $\card I \leqslant \pi$. On trouve alors $\Filt^{\pi}$, qui risque d'être une catégorie strictement plus petite que $\mathcal{M}_{\pi}$. Ce n'est pas non plus la catégorie des $\varinjlim\limits_{I} x_i$, avec $x_i \in C$ $\forall\ i \in \Ob I$, et $\card \Fl I \leqslant \pi$ -- puisque cette catégorie-là est $\Filt'^{\pi}$.)

\textsc{Démonstration}. Soit $\mathcal{M}'$ la catégorie de l'énoncé, on a donc $\mathcal{M}' \subset \mathcal{M}_{\pi}$, puisque $\mathcal{M}_{\pi}$ contient $C$ et est stable par les $\varinjlim$ de cardinal $\leqslant \pi$. Inversement, le lemme 4.5.14, $\mathrm{a} \Longrightarrow \mathrm{c}$.
% foot: Version 17 décembre 2025
\end{document}
